\section{AMC 12 12B 2021}
        \begin{enumerate}
        	\item (\textit{AMC 12 12B 2021}) How many integer values of $x$ satisfy $|x|<3\pi?$
$\textbf{(A) }9 \qquad \textbf{(B) }10 \qquad \textbf{(C) }18 \qquad \textbf{(D) }19 \qquad \textbf{(E) }20$

 \textbf{Answer: }D



\textbf{Solution 1}

Since $3\pi\approx9.42$, we multiply $9$ by $2$ for the integers from $1$ to $9$ and the integers from $-1$ to $-9$ and add $1$ to account for $0$ to get $\boxed{\textbf{(D)} ~19}$.

 



\textbf{Solution 2}

$|x|<3\pi$ $\iff$ $-3\pi<x<3\pi$. Since $\pi$ is approximately $3.14$, $3\pi$ is approximately $9.42$. We are trying to solve for $-9.42<x<9.42$, where $x\in\mathbb{Z}$. Hence, $-9.42<x<9.42$ $\implies$ $-9\leq x\leq9$, for $x\in\mathbb{Z}$. The number of integer values of $x$ is $9-(-9)+1=19$. Therefore, the answer is $\boxed{\textbf{(D)} ~19}$.

 



\textbf{Solution 3}

$3\pi \approx 9.4.$ There are two cases here.

 When $x>0, |x|>0,$ and $x = |x|.$ So then $x<9.4$

 When $x<0, |x|>0,$ and $x = -|x|.$ So then $-x<9.4$. Dividing by $-1$ and flipping the sign, we get $x>-9.4.$

 From case 1 and 2, we know that $-9.4 < x < 9.4$. Since $x$ is an integer, we must have $x$ between $-9$ and $9$. There are a total of 
 \[9-(-9) + 1 = \boxed{\textbf{(D)} ~19} \text{ integers}.\]~\href{https://artofproblemsolving.com/wiki/index.php?title=User:PureSwag}{PureSwag}

 



\textbf{Solution 4}

Looking at the problem, we see that instead of directly saying $x$, we see that it is $|x|.$ That means all the possible values of $x$ in this case are positive and negative. Rounding $\pi$ to $3$ we get $3(3)=9.$ There are $9$ positive solutions and $9$ negative solutions: $9+9=18.$ But what about zero? Even though zero is neither negative nor positive, but we still need to add it into the solution. Hence, the answer is $9+9+1=18+1=\boxed{\textbf{(D)} ~19}.$




\textbf{Solution 5}

There are an odd number of integer solutions $x$ to this inequality since if any non-zero integer $x$ satisfies this inequality, then so does $-x,$ and we must also account for $0,$ which gives us the desired. Then, the answer is either $\textbf{(A)}$ or $\textbf{(D)},$ and since $3 \pi > 3 \cdot 3 > 9,$ the answer is at least $9 \cdot 2 + 1 = 19,$ yielding $\boxed{\textbf{(D)} ~19}.$



	\item (\textit{AMC 12 12B 2021}) At a math contest, $57$ students are wearing blue shirts, and another $75$ students are wearing yellow shirts. The $132$ students are assigned into $66$ pairs. In exactly $23$ of these pairs, both students are wearing blue shirts. In how many pairs are both students wearing yellow shirts?


 $\textbf{(A) }23 \qquad \textbf{(B) }32 \qquad \textbf{(C) }37 \qquad \textbf{(D) }41 \qquad \textbf{(E) }64$

 \textbf{Answer: }B



\textbf{Solution}

There are $46$ blue students paired with a blue partner. The other $11$ students wearing blue shirts must each be paired with a partner wearing a shirt of the opposite color. There are $64$ yellow students remaining. Therefore the requested number of pairs is $\frac{64}{2}=\boxed{\textbf{(B)} ~32}$



	\item (\textit{AMC 12 12B 2021}) Suppose
 \[2+\frac{1}{1+\frac{1}{2+\frac{2}{3+x}}}=\frac{144}{53}.\]What is the value of $x?$
$\textbf{(A) }\frac34 \qquad \textbf{(B) }\frac78 \qquad \textbf{(C) }\frac{14}{15} \qquad \textbf{(D) }\frac{37}{38} \qquad \textbf{(E) }\frac{52}{53}$

 \textbf{Answer: }A



\textbf{Solution}

Subtracting $2$ from both sides and taking reciprocals gives $1+\frac{1}{2+\frac{2}{3+x}}=\frac{53}{38}$. Subtracting $1$ from both sides and taking reciprocals again gives $2+\frac{2}{3+x}=\frac{38}{15}$. Subtracting $2$ from both sides and taking reciprocals for the final time gives $\frac{x+3}{2}=\frac{15}{8}$ or $x=\frac{3}{4} \implies \boxed{\text{(A) } 3/4}$.

 


	\item (\textit{AMC 12 12B 2021}) Ms. Blackwell gives an exam to two classes. The mean of the scores of the students in the morning class is $84$, and the afternoon class's mean score is $70$. The ratio of the number of students in the morning class to the number of students in the afternoon class is $\frac34$. What is the mean of the score of all the students?


 $\textbf{(A) }74 \qquad \textbf{(B) }75 \qquad \textbf{(C) }76 \qquad \textbf{(D) }77 \qquad \textbf{(E) }78$

 \textbf{Answer: }C



\textbf{Solution 1 (One Variable)}

Let there be $3x$ students in the morning class and $4x$ students in the afternoon class. The total number of students is $3x + 4x = 7x$. The average is $\frac{3x\cdot84 + 4x\cdot70}{7x}=76$. Therefore, the answer is $\boxed{\textbf{(C)} ~76}$.

 



\textbf{Solution 2 (Two Variables)}

Suppose the morning class has $m$ students and the afternoon class has $a$ students. We have the following table:
 \[\begin{array}{c|c|c|c} & & & \\ [-2.5ex] & \textbf{\# of Students} & \textbf{Mean} & \textbf{Total} \\ \hline & & & \\ [-2.5ex] \textbf{Morning} & m & 84 & 84m \\ \hline & & & \\ [-2.5ex] \textbf{Afternoon} & a & 70 & 70a \end{array}\]We are also given that $\frac ma=\frac34,$ which rearranges as $m=\frac34a.$

 The mean of the scores of all the students is 
 \[\frac{84m+70a}{m+a}=\frac{84\left(\frac34a\right)+70a}{\frac34a+a}=\frac{133a}{\frac74a}=133\cdot\frac47=\boxed{\textbf{(C)} ~76}.\]~MRENTHUSIASM

 



\textbf{Solution 3 (Ratio)}

Of the average, $\frac{3}{3+4}=\frac{3}{7}$ of the scores came from the morning class and $\frac{4}{7}$ came from the afternoon class. The average is $\frac{3}{7}\cdot 84+\frac{4}{7}\cdot 70=\boxed{\textbf{(C)} ~76}.$




\textbf{Solution 4 (Convenient Values)}

WLOG, assume there are $3$ students in the morning class and $4$ in the afternoon class. Then the average is $\frac{3\cdot 84 + 4\cdot 70}{7}=\boxed{\textbf{(C)} ~76}.$




\textbf{Solution 5 (Basic manipulation)}

Let $S_1$ be the sum of the morning class's scores, and $S_2$ be the sum of the afternoon class's scores. Let $3x$ be the number of students in the morning class, and $4x$ be the number of students in the afternoon class. We can write $\frac{S_1}{3x}=84$ and $\frac{S_2}{4x}=70$, so $S_2 = 70 \cdot 4x$ and $S_1 = 84 \cdot 3x$. Adding these, we get $S_1 + S_2 = x(70 \cdot 4 + 84 \cdot 3)$. We want to find the mean, which is $\frac{S_1+S_2}{7x}$, so dividing by $7x$ on both sides of the equation, we get the mean to be $10 \cdot 4 + 12 \cdot 3 = \boxed{\textbf{(C)} ~76}.$



	\item (\textit{AMC 12 12B 2021}) The point $P(a,b)$ in the $xy$-plane is first rotated counterclockwise by $90^\circ$ around the point $(1,5)$ and then reflected about the line $y=-x$. The image of $P$ after these two transformations is at $(-6,3)$. What is $b-a?$
$\textbf{(A) }1 \qquad \textbf{(B) }3 \qquad \textbf{(C) }5 \qquad \textbf{(D) }7 \qquad \textbf{(E) }9$

 \textbf{Answer: }D



\textbf{Solution 1 (Transformation Rules)}

The final image of $P$ is $(-6,3)$. We know the reflection rule for reflecting over $y=-x$ is $(x,y) \rightarrow (-y, -x)$. So before the reflection and after rotation the point is $(-3,6)$.

 By definition of rotation, the slope between $(-3,6)$ and $(1,5)$ must be perpendicular to the slope between $(a,b)$ and $(1,5)$. The first slope is $\frac{5-6}{1-(-3)} = \frac{-1}{4}$. This means the slope of $P$ and $(1,5)$ is $4$.

 Rotations also preserve distance to the center of rotation, and since we only "travelled" up and down by the slope once to get from $(-3,6)$ to $(1,5)$ it follows we shall only use the slope once to travel from $(1,5)$ to $P$.

 Therefore point $P$ is located at $(1+1, 5+4) = (2,9)$. The answer is $9-2 = 7 = \boxed{\textbf{(D)} ~7}$.

 -abhinavg0627

 



\textbf{Solution 2 (Complex Numbers)}

Let us reconstruct that coordinate plane as the complex plane. Then, the point $P(a, b)$ becomes $a+b\cdot{i}$. A $90^\circ$ rotation around the point $(1, 5)$ can be done by translating the point $(1, 5)$ to the origin, rotating around the origin by $90^\circ$, and then translating the origin back to the point $(1, 5)$.
 \[a+b\cdot{i}  \implies (a-1)+(b-5)\cdot{i} \implies ((a-1)+(b-5)\cdot{i})\cdot{i} = 5-b+(a-1)i \implies 5+1-b+(a-1+5)i = 6-b+(a+4)i.\] By basis reflection rules, the reflection of $(-6, 3)$ about the line $y = -x$ is $(-3, 6)$. Hence, we have 
 \[6-b+(a+4)i = -3+6i \implies b=9, a=2,\] from which $b-a = 9-2 = \boxed{\textbf{(D)} ~7}$. 

 



\textbf{Solution 3 (Reverso)}

The problem gives a series of transformations and proceeds to give the resultant point, being $(-6,3)=P$. Therefore, all we must do is reverse it. First, we reverse the last transformation by computing the distance from the point $(-6,3)$ to $y=-x$ by using the formula $d=\frac{|Ax+By+C|}{\sqrt{A^2+B^2}}$. Where $Ax+By+C$ is the standard form of a line. Computing, we obtain that $d=\frac{3\sqrt{2}}{2}$. We also know that this is magnitude is at an angle of $\frac{\pi}{4}$. Therefore, to do the transformation, we double that vector and add it to the point. We get:

 $P^{'}=P+2\overrightarrow{v}$$\implies P'=(-6,3)+\left\langle 3\sqrt{2}\cdot \frac{\sqrt{2}}{2},3\sqrt{2}\cdot \frac{\sqrt{2}}{2} \right \rangle$$\implies P'=(-3,6)$. 

 Now, we must reverse the second transformation. To do so, realize that $P'-(1,5)\Leftrightarrow \overrightarrow{v_{2}}=\left\langle -4,1 \right\rangle$. Simply make this vector perpendicular by switching the $x$ and $y$ components and switching the sign of the initial $y$ component. Therefore, we get $\overrightarrow{v_{2\bot }}=\left\langle 1,4 \right\rangle$. Therefore, adding this vector to $(1,5)$ yields $(2,9)$, which leads us to $9-2 = \boxed{\textbf{(D)} ~7}$.

 



\textbf{Solution 4 (Vector Dot Product)}

Using the same method as in Solution 1, we can obtain that the point before the reflection is $(-3,6)$. If we let the original point be $(x, y)$, then we can use that the starting point is $(1,5)$ to obtain two vectors $\langle -4,1 \rangle$ and $\langle x-1, y-5 \rangle$.We know that two vectors are perpendicular if their dot product is equal to $0$, and that both points are the same distance ($\sqrt {17}$) from $(1,5)$. 

 Therefore, we can write two equations using these vectors:$(x-1)^2 + (y-5)^2 = 17$ (from distance and pythagorean theorem) and $-4x+y-1 = 0$ (from dot product)

 Solving, we simplify the second equation to $y=4x+1$, and plug it into the first equation.We obtain $(x-1)^2 + (4x-4)^2 = 17$. We can simplify this to the quadratic $17x^2-34x=0$. When we factor out $17x$, we find that $x = 2$ or $x = 0$. However, $x$ cannot equal $0$.Therefore, $x = 2$, and plugging this into the second equation gives us that $y = 9$. Since the point is $(9, 2)$, we compute $9-2 = \boxed{\textbf{(D)} ~7}$.

 


 



\textbf{Solution 5 (Vector Dot Product scuffed version)}

Using the same method as in Solution 1 reflecting $(-6,3)$ about the line $y = -x$ gives us $(-3,6).$

 Let the original point be $\langle x,y \rangle.$ From point $(1,5),$ we form the vectors $\langle -4,1 \rangle$ and $\langle x-1, y-5 \rangle$ that extend out from the initial point. If they are perpendicular, we know that their dot product has to equal zero. Therefore, 
 \[\langle -4,1 \rangle \cdot \langle x-1, y-5 \rangle = 0 \implies -4x+y-1= 0.\]Now, we have to do some guess and check from the multiple choices. Let $y - x = A$ where $A$ is one of the answer choices. Then, $A -3x = 1.$ By intuition and logical reasoning we deduce that $A$ must be $1 \pmod 3$ so that brings our potential answers down to $\text{\textbf{(A)}}$ and $\text{\textbf{(C)}}.$ If $A = 1$ from $\text{\textbf{(A)}},$ then $x = 0,$ which we can quickly rule out since we know thar $P$ rotated counterclockwise not clockwise. Hence, $\boxed{\textbf{(D)} ~7}$ is the answer.

 


 



\textbf{Solution 6 (Working Backwards Approach, Super Easy!)}

The last thing that happened was the reflection over the line $y = -x$, so we reverse that by reflecting $(-6,3)$ over $y = -x$ to get $(-3,6)$. Now, we have to rotate it $\textbf{clockwise}$, to reverse the first step. To rotate $(-3,6)$ over the point $(1,5)$ clockwise 90 degrees, we can subtract $1$ from each x-coordinate and subtract $5$ from each y-coordinate. That way, we only have to rotate $(-4,1)$ clockwise 90 degrees over the origin $(0,0)$. This is pretty easy to do in your head, so you get $(1,4)$. But we are not, done, we still have to add $1$ to each x-coordinate and add $5$ to each y-coordinate, to get $(2,9)$, so our final answer is $9 - 2 = 7$, which is $\boxed{\textbf{(D)} ~7}$.

 


	\item (\textit{AMC 12 12B 2021}) An inverted cone with base radius $12 \text{cm}$ and height $18\text{cm}$ is full of water. The water is poured into a tall cylinder whose horizontal base has a radius of $24\text{cm}$. What is the height in centimeters of the water in the cylinder?


 $\textbf{(A) }1.5 \qquad \textbf{(B) }3 \qquad \textbf{(C) }4 \qquad \textbf{(D) }4.5 \qquad \textbf{(E) }6$

 \textbf{Answer: }A



\textbf{Solution 1}

The volume of a cone is $\frac{1}{3} \cdot\pi \cdot r^2 \cdot h$ where $r$ is the base radius and $h$ is the height. The water completely fills up the cone so the volume of the water is $\frac{1}{3}\cdot18\cdot144\pi = 6\cdot144\pi$.

 The volume of a cylinder is $\pi \cdot r^2 \cdot h$ so the volume of the water in the cylinder would be $24\cdot24\cdot\pi\cdot h$.

 We can equate these two expressions because the water volume stays the same like this $24\cdot24\cdot\pi\cdot h = 6\cdot144\pi$. We get $4h = 6$ and $h=\frac{6}{4}$.

 So the answer is $\boxed{\textbf{(A)} ~1.5}.$




\textbf{Solution 2}

The water completely fills up the cone. For now, assume the radius of both cone and cylinder are the same. Then the cone has $\frac{1}{3}$ of the volume of the cylinder, and so the height is divided by $3$. Then, from the problem statement, the radius is doubled, meaning the area of the base is quadrupled (since $2^2 = 4$).

 Therefore, the height is divided by $3$ and divided by $4$, which is $18 \div 3 \div 4 = \boxed{\textbf{(A)} ~1.5}.$



	\item (\textit{AMC 12 12B 2021}) Let $N=34\cdot34\cdot63\cdot270.$ What is the ratio of the sum of the odd divisors of $N$ to the sum of the even divisors of $N?$
$\textbf{(A) }1:16 \qquad \textbf{(B) }1:15 \qquad \textbf{(C) }1:14 \qquad \textbf{(D) }1:8 \qquad \textbf{(E) }1:3$

 \textbf{Answer: }C



\textbf{Solution 1}

Prime factorize $N$ to get $N=2^{3} \cdot 3^{5} \cdot 5\cdot 7\cdot 17^{2}$. For each odd divisor $n$ of $N$, there exist even divisors $2n, 4n, 8n$ of $N$, therefore the ratio is $1:(2+4+8)=\boxed{\textbf{(C)} ~1 : 14}$




\textbf{Solution 2}

Prime factorizing $N$, we see $N=2^{3} \cdot 3^{5} \cdot 5\cdot 7\cdot 17^{2}$.The sum of $N$'s odd divisors are the sum of the factors of $N$ without $2$, and the sum of the even divisors is the sum of the odds subtracted by the total sum of divisors. 

 BY SUM OF FACTORS FORMULA (search it up if you don't know): The formula is actually for all factors, but we can just take out the $2^3$, so we have:

 The sum of odd divisors is given by 
 \[a = (1+3+3^2 + 3^3 + 3^4 + 3^5)(1 + 5)(1+7)(1+17+17^2)\] and the total sum of divisors is 
 \[(1+2+4+8)(1+3+3^2 + 3^3 + 3^4 + 3^5)(1 + 5)(1+7)(1+17+17^2) = 15a.\] Thus, our ratio is 
 \[\frac{a}{15a-a} = \frac{a}{14a} = \boxed{\textbf{(C)} ~1 : 14}.\]




\textbf{Solution 3}

Prime factorizing $N$, we have that there is $2^3$ in our factorization. Now, call the sum of the odd divisors $k$. We know that if we multiply k by 2, we will have even divisors. So, we can multiply k by 2, $2^2, 2^3$ respectively to get 14k as the sum of the even divisors. Therefore, the answer is 
 \[\frac{k}{14k} = \boxed{\textbf{(C)} ~1:14}.\]~MC

 


	\item (\textit{AMC 12 12B 2021}) Three equally spaced parallel lines intersect a circle, creating three chords of lengths $38,38,$ and $34$. What is the distance between two adjacent parallel lines?


 $\textbf{(A) }5\frac12 \qquad \textbf{(B) }6 \qquad \textbf{(C) }6\frac12 \qquad \textbf{(D) }7 \qquad \textbf{(E) }7\frac12$

 \textbf{Answer: }B



\textbf{Solution 1 (Pythagorean Theorem)}


                    \begin{figure}[H]
                        \centering
                        \includegraphics[width=0.5\textwidth]{Solution/AMC_12_12B_2021_8_Solution_Figure_1.png}
                    \end{figure}
                    

 


 Since two parallel chords have the same length ($38$), they must be equidistant from the center of the circle. Let the perpendicular distance of each chord from the center of the circle be $d$. Thus, the distance from the center of the circle to the chord of length $34$ is 

 
 \[2d + d = 3d\]

 and the distance between each of the chords is just $2d$. Let the radius of the circle be $r$. Drawing radii to the points where the lines intersect the circle, we create two different right triangles: 

 - One with base $\frac{38}{2}= 19$, height $d$, and hypotenuse $r$  ($\triangle RAO$ on the diagram)

 - Another with base $\frac{34}{2} = 17$, height $3d$, and hypotenuse $r$  ($\triangle LBO$ on the diagram)

 By the Pythagorean theorem, we can create the following system of equations:

 
 \[19^2 + d^2 = r^2\]

 \[17^2 + (2d + d)^2 = r^2\]

 Solving, we find $d = 3$, so $2d = \boxed{\textbf{(B)}\ 6}$.

 



\textbf{Solution 2 (Coordinates)}

Because we know that the equation of a circle is $(x-a)^2 + (y-b)^2 = r^2$ where the center of the circle is $(a, b)$ and the radius is $r$, we can find the equation of this circle by centering it on the origin. Doing this, we get that the equation is $x^2 + y^2 = r^2$. Now, we can set the distance between the chords as $2d$ so the distance from the chord with length $38$ to the diameter is $d$. 

 Therefore, the following points are on the circle as the y-axis splits the chord in half, that is where we get our $x$ value:

 $(19, d)$
$(19, -d)$
$(17, -3d)$

 
Now, we can plug one of the first two value in as well as the last one to get the following equations:

 
 \[19^2 + d^2 = r^2\]

 \[17^2 + (3d)^2 = r^2\]

 Subtracting these two equations, we get $19^2 - 17^2 = 8d^2$ - therefore, we get $72 = 8d^2 \rightarrow d^2 = 9 \rightarrow d = 3$. We want to find $2d = 6$ because that's the distance between two chords. So, our answer is $\boxed{B}$.

 



\textbf{Solution 3 (Stewart's Theorem)}


                    \begin{figure}[H]
                        \centering
                        \includegraphics[width=0.5\textwidth]{Solution/AMC_12_12B_2021_8_Solution_Figure_2.png}
                    \end{figure}
                    If $d$ is the requested distance, and $r$ is the radius of the circle, Stewart's Theorem applied to $\triangle OCD$ with cevian $\overleftrightarrow{OP}$ gives 
 \[19\cdot 38\cdot 19 + \frac{1}{2}d\cdot 38\cdot\frac{1}{2}d=19r^{2}+19r^{2}.\] This simplifies to $13718+\frac{19}{2}d^{2}=38r^{2}$. Similarly, another round of Stewart's Theorem applied to $\triangle OEF$ with cevian $\overleftrightarrow{OQ}$ gives 
 \[17\cdot 34\cdot 17 + \frac{3}{2}d\cdot 34\cdot\frac{3}{2}d=17r^{2}+17r^{2}.\] This simplifies to $9826+\frac{153}{2}d^{2}=34r^{2}$. Dividing the top equation by $38$ and the bottom equation by $34$ results in the system of equations 
 \begin{align*} 361+\frac{1}{4}d^{2} &= r^{2} \\ 289+\frac{9}{4}d^{2} &= r^{2} \\ \end{align*}
By transitive, $361+\frac{1}{4}d^{2}=289+\frac{9}{4}d^{2}$. Therefore $(\frac{9}{4}-\frac{1}{4})d^{2}=361-289\rightarrow 2d^{2}=72\rightarrow d^{2}=36\rightarrow d=\boxed{\textbf{(B)} ~6}.$



	\item (\textit{AMC 12 12B 2021}) What is the value of
 \[\frac{\log_2 80}{\log_{40}2}-\frac{\log_2 160}{\log_{20}2}?\]
$\textbf{(A) }0 \qquad \textbf{(B) }1 \qquad \textbf{(C) }\frac54 \qquad \textbf{(D) }2 \qquad \textbf{(E) }\log_2 5$

 \textbf{Answer: }D



\textbf{Solution}


 \[\frac{\log_{2}{80}}{\log_{40}{2}}-\frac{\log_{2}{160}}{\log_{20}{2}}\]

 Note that $\log_{40}{2}=\frac{1}{\log_{2}{40}}$, and similarly $\log_{20}{2}=\frac{1}{\log_{2}{20}}$

 \[= \log_{2}{80}\cdot \log_{2}{40}-\log_{2}{160}\cdot \log_{2}{20}\]

 \[=(\log_{2}{4}+\log_{2}{20})(\log_{2}{2}+\log_{2}{20})-(\log_{2}{8}+\log_{2}{20})\log_{2}{20}\]

 \[=(2+\log_{2}{20})(1+\log_{2}{20})-(3+\log_{2}{20})\log_{2}{20}\]

 Expanding, 
 \[2+2\log_{2}{20}+\log_{2}{20}+(\log_{2}{20})^2-3\log_{2}{20}-(\log_{2}{20})^2\]

 All the log terms cancel, so the answer is $2\implies\boxed{\text{(D)}}$.

 


	\item (\textit{AMC 12 12B 2021}) Two distinct numbers are selected from the set $\{1,2,3,4,\dots,36,37\}$ so that the sum of the remaining $35$ numbers is the product of these two numbers. What is the difference of these two numbers?


 $\textbf{(A) }5 \qquad \textbf{(B) }7 \qquad \textbf{(C) }8\qquad \textbf{(D) }9 \qquad \textbf{(E) }10$

 \textbf{Answer: }E



\textbf{Solution}

The sum of the first $n$ integers is given by $\frac{n(n+1)}{2}$, so $\frac{37(37+1)}{2}=703$.

 Therefore, $703-x-y=xy$

 Rearranging, $xy+x+y=703$. We can factor this equation by \href{https://artofproblemsolving.com/wiki/index.php?title=SFFT}{SFFT} to get

 $(x+1)(y+1)=704$

 Looking at the possible divisors of $704 = 2^6\cdot11$, $22$ and $32$ are within the constraints of $0 < x \leq y \leq 37$ so we try those:

 $(x+1)(y+1) = 22\cdot32$
$x+1=22, y+1 = 32$
$x = 21, y = 31$

 Therefore, the difference $y-x=31-21=\boxed{\textbf{(E) }10}$.

 


 


	\item (\textit{AMC 12 12B 2021}) Triangle $ABC$ has $AB=13,BC=14$ and $AC=15$. Let $P$ be the point on $\overline{AC}$ such that $PC=10$. There are exactly two points $D$ and $E$ on line $BP$ such that quadrilaterals $ABCD$ and $ABCE$ are trapezoids. What is the distance $DE?$
$\textbf{(A) }\frac{42}5 \qquad \textbf{(B) }6\sqrt2 \qquad \textbf{(C) }\frac{84}5\qquad \textbf{(D) }12\sqrt2 \qquad \textbf{(E) }18$

 \textbf{Answer: }D



\textbf{Solution 1 (Analytic Geometry)}

Toss on the Cartesian plane with $A=(5, 12), B=(0, 0),$ and $C=(14, 0)$. Then $\overline{AD}\parallel\overline{BC}, \overline{AB}\parallel\overline{CE}$ by the trapezoid condition, where $D, E\in\overline{BP}$. Since $PC=10$, point $P$ is $\frac{10}{15}=\frac{2}{3}$ of the way from $C$ to $A$ and is located at $(8, 8)$. Thus line $BP$ has equation $y=x$. Since $\overline{AD}\parallel\overline{BC}$ and $\overline{BC}$ is parallel to the ground, we know $D$ has the same $y$-coordinate as $A$, except it'll also lie on the line $y=x$. Therefore, $D=(12, 12). \, \blacksquare$

 To find the location of point $E$, we need to find the intersection of $y=x$ with a line parallel to $\overline{AB}$ passing through $C$. The slope of this line is the same as the slope of $\overline{AB}$, or $\frac{12}{5}$, and has equation $y=\frac{12}{5}x-\frac{168}{5}$. The intersection of this line with $y=x$ is $(24, 24)$. Therefore point $E$ is located at $(24, 24). \, \blacksquare$

 The distance $DE$ is equal to the distance between $(12, 12)$ and $(24, 24)$, which is $\boxed{\textbf{(D) }12\sqrt2}$.

 



\textbf{Solution 2 (Stewart's Theorem)}

Using Stewart's Theorem we find $BP = 8\sqrt{2}$. From the similar triangles $BPA\sim EPC$ and $BPC\sim DPA$, we have
 \[DP = BP\cdot\frac{PA}{PC} = \frac{1}{2}BP\]
 \[EP = BP\cdot\frac{PC}{PA} = 2BP\]So, 
 \[DE = \frac{3}{2}BP = \boxed{\textbf{(D) }12\sqrt2}.\]




\textbf{Solution 3}

Let $x$ be the length $PE$. From the similar triangles $BPA\sim EPC$ and $BPC\sim DPA$ we have
 \[BP = \frac{PA}{PC}x = \frac12 x\]
 \[PD = \frac{PA}{PC}BP = \frac14 x\]Therefore $BD = DE = \frac{3}{4}x$. Now extend line $CD$ to the point $Z$ on $AE$, forming parallelogram $ZABC$. As $BD = DE$ we also have $EZ = ZC = 13$ so $EC = 26$.

 We now use the Law of Cosines to find $x$ (the length of $PE$):
 \[x^2 = EC^2 + PC^2 - 2(EC)(PC)\cos{(PCE)} = 26^2 + 10^2 - 2\cdot 26\cdot 10\cos(\angle PCE)\]As $\angle PCE = \angle BAC$, we have (by Law of Cosines on triangle $BAC$)
 \[\cos(\angle PCE) = \frac{13^2 + 15^2 - 14^2}{2\cdot 13\cdot 15}.\]Therefore
 \begin{align*} x^2 &= 26^2 + 10^2 - 2\cdot 26\cdot 10\cdot\frac{198}{2\cdot 13\cdot 15}\\ &= 776 - 264\\ &= 512 \end{align*}
And $x = 16\sqrt2$. The answer is then $\frac34x = \boxed{\textbf{(D) }12\sqrt2}$.

 



\textbf{Solution 4 (Heron's Formula, Pythagorean Theorem, Similar Triangles)}

Let the brackets denote areas. By Heron's Formula, we have
 \begin{align*} [ABC]&=\sqrt{\frac{13+14+15}{2}\left(\frac{13+14+15}{2}-13\right)\left(\frac{13+14+15}{2}-14\right)\left(\frac{13+14+15}{2}-15\right)} \\ &=\sqrt{21\left(21-13\right)\left(21-14\right)\left(21-15\right)} \\ &=\sqrt{21\left(8\right)\left(7\right)\left(6\right)} \\ &=\sqrt{\left(3\cdot7\right)\left(2^3\right)\left(7\right)\left(2\cdot3\right)} \\ &=2^2\cdot3\cdot7 \\ &=84. \end{align*}
It follows that the height of $ABCD$ is $\frac{2[ABC]}{14}=12.$

 Next, we drop the altitudes $\overline{AF}$ and $\overline{DG}$ of $ABCD.$ By the Pythagorean Theorem on right $\triangle AFB,$ we get $BF=5.$ By the AA Similarity, we have $\triangle ADP\sim\triangle CBP,$ with the ratio of similitude $1:2.$ It follows that $AD=7.$ Since $ADGF$ is a rectangle, we get $FG=AD=7.$ By the Pythagorean Theorem on right $\triangle DGB,$ we get $BD=12\sqrt2.$

 By $\triangle ADP\sim\triangle CBP$ again, we get $BP=8\sqrt2$ and $DP=4\sqrt2.$ Also, by the AA Similarity, we have $\triangle ABP\sim\triangle CEP,$ with the ratio of similitude $1:2.$ It follows that $EP=16\sqrt2.$

 Finally, we obtain $DE=EP-DP=\boxed{\textbf{(D) }12\sqrt2}.$
<u>\textbf{Remark}</u>

 It is well known that the area of a $13\text{-}14\text{-}15$ triangle is $84.$




\textbf{Solution 5 (Barycentric Coordinates)}

(For those unfamiliar with barycentric coordinates, consider reading the barycentric coordinates article written by Evan Chen and Max Schindler here: \href{https://artofproblemsolving.com/wiki/index.php/Resources\_for\_mathematics\_competitions\#Articles}{https://artofproblemsolving.com/wiki/index.php/Resources\_for\_mathematics\_competitions\#Articles})

 We can find $P$ in barycentric coordinates as $\Bigr(\frac{2}{3},0,\frac{1}{3}\Bigr)$. We can then write $\overline{BP}$ as $x-2z=0$, where $(x,y,z)$ defines a point in barycentric coordinates. We have $\overline{AD}\parallel \overline{BC}$ as $y+z=0$ and $\overline{CE}\parallel \overline{AB}$ as $x+y=0$. We can then compute $D$ and $E$ by intersecting lines:
 \[\begin{cases} x-2z=0\\ y+z=0\\ x+y+z=1 \end{cases}\]Which gives us $D=\left(1, -\frac{1}{2}, \frac{1}{2}\right)$. We can get $E$ with:
 \[\begin{cases} x-2z=0\\ x+y=0\\ x+y+z=1 \end{cases}\]Which gives us $E=(2, -2, 1)$. Then, finding the displacement vector, we have $\overrightarrow{ED}=\left(1,-\frac{3}{2},\frac{1}{2}\right)$. Using the barycentric distance formula:
 \begin{align*} \text{dist}(D,E)&=\sqrt{-a^2yx-b^2zx-c^2xy} \\ &=\sqrt{-14^2\Bigr(-\frac{3}{2}\Bigr)\Bigr(\frac{1}{2}\Bigr)-15^2\Bigr(1\Bigr)\Bigr(\frac{1}{2}\Bigr)-13^2\Bigr(-\frac{3}{2}\Bigr)\Bigr(\frac{1}{2}\Bigr)} \\ &=\boxed{\textbf{(D) }12\sqrt2}. \end{align*}




	\item (\textit{AMC 12 12B 2021}) Suppose that $S$ is a finite set of positive integers. If the greatest integer in $S$ is removed from $S$, then the average value (arithmetic mean) of the integers remaining is $32$. If the least integer in $S$ is also removed, then the average value of the integers remaining is $35$. If the greatest integer is then returned to the set, the average value of the integers rises to $40$. The greatest integer in the original set $S$ is $72$ greater than the least integer in $S$. What is the average value of all the integers in the set $S$?


 $\textbf{(A) }36.2 \qquad \textbf{(B) }36.4 \qquad \textbf{(C) }36.6\qquad \textbf{(D) }36.8 \qquad \textbf{(E) }37$

 \textbf{Answer: }D



\textbf{Solution 1}

We can then say that \( \frac{A+S(n)}{n+1} = 32 \), \( \frac{S(n)}{n} = 35 \), and \( \frac{B+S(n)}{n+1} = 40 \). 

 Expanding gives us \( A+S(n) = 32n+32 \), \( S(n) = 35n \), and \( B+S(n) = 40n+40 \). 

 Substituting \( S(n) = 35n \) to all gives us \( A+35n=32n+32 \) and \( B+35n=40n+40 \). 

 Solving for \( A \) and \( B \) gives \( A=-3n+32 \) and \( B = 5n+40 \). 

 We now need to find \( \frac{S(n)+A+B}{n+2} \). We substitute everything to get \( \frac{35n+(-3n+32)+(5n+40)}{n+2} \), or \( \frac{37n+72}{n+2} \). 

 Say that the answer to this is \( Z \). Then, \( Z \) needs to be a number that makes \( n \) a positive integer. The only options that work is $\boxed{\textbf{(C) }36.6}$ and $\boxed{\textbf{(D) }36.8}$. 

 However, if 36.6 is an option, we get \( n=3 \). So that means that \(A\) is \(23\) and \(B\) is \(55\), and \( S(n)=105 \). But if there is \(3\) terms, then the middle number is \(105\), but we said that \( B \) is the largest number in the set, so therefore our answer cannot be $\boxed{\textbf{(C) }36.6}$ and is instead $\boxed{\textbf{(D) }36.8}$ and now, we're finished!

 



\textbf{Solution 1 Alternative Ending}

If A is smaller than B by 72 therefore from the equation on the top you can find out that \(N = 8\) using substitution the plug it in to the equation \( \frac{37n+72}{n+2} \) then you will get that Z = $\boxed{\textbf{(D) }36.8}$. 

 This was more efficient then the last part.

 



\textbf{Solution 2}

Let the lowest value be $L$ and the highest $G$, and let the sum be $Z$ and the amount of numbers $n$. We have $\frac{Z-G}{n-1}=32$, $\frac{Z-L-G}{n-2}=35$, $\frac{Z-L}{n-1}=40$, and $G=L+72$. Clearing denominators gives $Z-G=32n-32$, $Z-L-G=35n-70$, and $Z-L=40n-40$. We use $G=L+72$ to turn the first equation into $Z-L=32n+40$. Since $Z-L=40(n-1)$ we substitute it into the equation which gives $n=10$. Turning the second into $Z-2L=35n+2$ using $G=L+72$ we see $L=8$ and $Z=368$ so the average is $\frac{Z}{n}=\boxed{\textbf{(D) }36.8}$ ~aop2014

 



\textbf{Solution 3}

Let $x$ be the greatest integer, $y$ be the smallest, $z$ be the sum of the numbers in S excluding $x$ and $y$, and $k$ be the number of elements in S.

 Then, $S=x+y+z$

 First, when the greatest integer is removed, $\frac{S-x}{k-1}=32$

 When the smallest integer is also removed, $\frac{S-x-y}{k-2}=35$

 When the greatest integer is added back, $\frac{S-y}{k-1}=40$

 We are given that $x=y+72$

 
After you substitute $x=y+72$, you have 3 equations with 3 unknowns $S,$, $y$ and $k$.

 $S-y-72=32k-32$
$S-2y-72=35k-70$
$S-y=40k-40$

 This can be easily solved to yield $k=10$, $y=8$, $S=368$.

 $\therefore$ average value of all integers in the set $=S/k = 368/10 = \boxed{\textbf{(D) }36.8}$




\textbf{Solution 4}

We should plug in $36.2$ and assume everything is true except the $35$ part. We then calculate that part and end up with $35.75$. We also see with the formulas we used with the plug in that when you increase by $0.2$ the $35.75$ part decreases by $0.25$. The answer is then $\boxed{\textbf{(D) }36.8}$. You can work backwards because it is multiple choice and you don't have to do critical thinking. ~Lopkiloinm

 



\textbf{Solution 5}

Let $S = \{a_1, a_2, a_3, \hdots, a_n\}$ with $a_1 < a_2 < a_3 < \hdots < a_n.$ We are given the following: 
 \[{\begin{cases} \sum_{i=1}^{n-1} a_i = 32(n-1) = 32n-32, \\ \sum_{i=2}^n a_i = 40(n-1) = 40n-40, \\ \sum_{i=2}^{n-1} a_i = 35(n-2) = 35n-70, \\ a_n-a_1 = 72 \implies a_1 + 72 = a_n. \end{cases}}\] Subtracting the third equation from the sum of the first two, we find that 
 \[\sum_{i=1}^n a_i = \left(32n-32\right) + \left(40n-40\right) - \left(35n-70\right) = 37n - 2.\] Furthermore, from the fourth equation, we have 
 \[\sum_{i=2}^{n} a_i - \sum_{i=1}^{n-1} a_i = \left[\left(a_1 + 72\right) + \sum_{i=2}^{n-1} a_i\right] - \left[\left(a_1\right) + \sum_{i=2}^{n-1} a_i\right] = \left(40n-40\right)-\left(32n-32\right).\] Combining like terms and simplifying, we have 
 \[72 = 8n-8 \implies 8n = 80 \implies n=10.\] Thus, the sum of the elements in $S$ is $37 \cdot 10 - 2 = 368,$ and since there are 10 elements in $S,$ the average of the elements in $S$ is $\frac{368}{10}=\boxed{\textbf{(D) }36.8}$




\textbf{Solution 6}

Let $n$ be the number of elements in $S, m_l = 32, \Sigma_l = m_l \cdot (n-1), m_g = 40,  \Sigma_g = m_g \cdot (n-1), m_{lg} = 35,  \Sigma_lg = m_{lg} \cdot (n-2).$
 \[\Sigma_g - \Sigma_l = 72 = (n-1) \cdot (m_g - m_l) \implies n = 1 + \frac {72}{m_g - m_l} = 1 + \frac {72}{40-32}=10.\]
 \[m = \frac {\Sigma_g + \Sigma_l - \Sigma_{lg}}{n} = \frac {40 \cdot 9 + 32 \cdot 9 - 35 \cdot 8}{10} = 36.8.\]\textbf{vladimir.shelomovskii@gmail.com, vvsss}

 


	\item (\textit{AMC 12 12B 2021}) How many values of $\theta$ in the interval $0<\theta\le 2\pi$ satisfy
 \[1-3\sin\theta+5\cos3\theta = 0?\]
$\textbf{(A) }2 \qquad \textbf{(B) }4 \qquad \textbf{(C) }5\qquad \textbf{(D) }6 \qquad \textbf{(E) }8$

 \textbf{Answer: }D



\textbf{Solution}

We rearrange to get 
 \[5\cos3\theta = 3\sin\theta-1.\]We can graph two functions in this case: $y=5\cos{3x}$ and $y=3\sin{x} -1$.Using transformation of functions, we know that $5\cos{3x}$ is just a cosine function with amplitude $5$ and period $\frac{2\pi}{3}$. Similarly, $3\sin{x} -1$ is just a sine function with amplitude $3$ and shifted $1$ unit downward:
                    \begin{figure}[H]
                        \centering
                        \includegraphics[width=0.5\textwidth]{Solution/AMC_12_12B_2021_13_Solution_Figure_1.png}
                    \end{figure}
                    So, we have $\boxed{\textbf{(D) }6}$ solutions. 

 


	\item (\textit{AMC 12 12B 2021}) Let $ABCD$ be a rectangle and let $\overline{DM}$ be a segment perpendicular to the plane of $ABCD$. Suppose that $\overline{DM}$ has integer length, and the lengths of $\overline{MA},\overline{MC},$ and $\overline{MB}$ are consecutive odd positive integers (in this order). What is the volume of pyramid $MABCD?$
$\textbf{(A) }24\sqrt5 \qquad \textbf{(B) }60 \qquad \textbf{(C) }28\sqrt5\qquad \textbf{(D) }66 \qquad \textbf{(E) }8\sqrt{70}$

 \textbf{Answer: }A



\textbf{Solution 1}

Let $MA=a$ and $MD=d.$ It follows that $MC=a+2$ and $MB=a+4.$

 As shown below, note that $\triangle MAD$ and $\triangle MBC$ are both right triangles.
                    \begin{figure}[H]
                        \centering
                        \includegraphics[width=0.5\textwidth]{Solution/AMC_12_12B_2021_14_Solution_Figure_1.png}
                    \end{figure}
                    By the Pythagorean Theorem, we have
 \begin{alignat*}{6} AD\textasciicircum{}2 &= MA\textasciicircum{}2 - MD\textasciicircum{}2 &&= a\textasciicircum{}2 - d\textasciicircum{}2, \\ BC\textasciicircum{}2 &= MB\textasciicircum{}2 - MC\textasciicircum{}2 &&= (a+4)\textasciicircum{}2 - (a+2)\textasciicircum{}2. \end{alignat*}Since $AD=BC$ in rectangle $ABCD,$ we equate the expressions for $AD^2$ and $BC^2,$ then rearrange and factor:
 \begin{align*} a^2 - d^2 &= (a+4)^2 - (a+2)^2 \\ a^2 - d^2 &= 4a + 12 \\ a^2 - 4a - d^2 &= 12 \\ (a-2)^2 - d^2 &= 16 \\ (a+d-2)(a-d-2) &= 16. \end{align*}
As $a+d-2$ and $a-d-2$ have the same parity, we get $a+d-2=8$ and $a-d-2=2,$ from which $(a,d)=(7,3).$

 Applying the Pythagorean Theorem to right $\triangle MAD$ and right $\triangle MCD,$ we obtain $AD=2\sqrt{10}$ and $CD=6\sqrt2,$ respectively.

 Let the brackets denote areas. Together, the volume of pyramid $MABCD$ is 
 \[\frac13\cdot [ABCD]\cdot MD = \frac13\cdot (AD\cdot CD)\cdot MD = \boxed{\textbf{(A) }24\sqrt5}.\]~Lopkiloinm ~MRENTHUSIASM

 



\textbf{Solution 2}

Let $AD=b$, $CD=a$, $MD=x$, $MC=t$. It follows that $MA=t-2$ and $MB=t+2$.

 We have three equations:
 \begin{align*} a^2 + x^2 &= t^2, \\ a^2 + b^2 + x^2 &= t^2 + 4t + 4, \\ b^2 + x^2 &= t^2 - 4t + 4. \end{align*}
Substituting the first and third equations into the second equation, we get:
 \begin{align*} t^2 - 8t - x^2 &= 0 \\ (t-4)^2 - x^2 &= 16 \\ (t-4-x)(t-4+x) &= 16. \end{align*}
Therefore, we have $t = 9$ and $x = 3$.

 Solving for other values, we get $b = 2\sqrt{10}$, $a = 6\sqrt{2}$.The volume is then 
 \[\frac{1}{3} abx = \boxed{\textbf{(A) }24\sqrt5}.\]




\textbf{Solution 3}

First and foremost, the diagram must be drawn.
                    \begin{figure}[H]
                        \centering
                        \includegraphics[width=0.5\textwidth]{Solution/AMC_12_12B_2021_14_Solution_Figure_2.png}
                    \end{figure}
                    With the diagram, an impulse to utilize Pythagorean theorem is created. Thereby, the shaded right triangles may be used.
                    \begin{figure}[H]
                        \centering
                        \includegraphics[width=0.5\textwidth]{Solution/AMC_12_12B_2021_14_Solution_Figure_3.png}
                    \end{figure}
                    

 The following equations could be obtained from Pythagorean Theorem.
 \begin{align*}     I^2+a^2&=x^2 \\     I^2+b^2&=(x+2)^2 \\     I^2+a^2+b^2&=(x+4)^2 \\ \end{align*}
Because $x$ and $I$ are integers, writing an equation in terms of only $x$ and $I$ may be beneficial for trial and error.
 \begin{align*}     x^2+(x+2)^2&=(x+4)^2+I^2 \\     2x^2+4x+4&=x^2+8x+16+I^2 \\     x^2-4x-12&=I^2 \\     (x-6)(x+2)&=I^2 \end{align*}
$x=7$ seems valid since $I=3$ when $x=7$. Therefore, when $x=7$, $a=2\sqrt{10}$ and $b=6\sqrt{2}$. We could double check that $40+72+9=11^2$. Thereby,
 \[\frac{a\cdot b\cdot I}{3}=\frac{2\sqrt{10}\cdot6\sqrt{2}\cdot3}{3}=\boxed{\textbf{(A) }24\sqrt5}.\]



	\item (\textit{AMC 12 12B 2021}) The figure is constructed from $11$ line segments, each of which has length $2$. The area of pentagon $ABCDE$ can be written as $\sqrt{m} + \sqrt{n}$, where $m$ and $n$ are positive integers. What is $m + n ?$

                    \begin{figure}[H]
                        \centering
                        \includegraphics[width=0.5\textwidth]{Problem/AMC_12_12B_2021_15_Problem_Figure_1.png}
                    \end{figure}
                    
$\textbf{(A)} ~20 \qquad\textbf{(B)} ~21 \qquad\textbf{(C)} ~22 \qquad\textbf{(D)} ~23 \qquad\textbf{(E)} ~24$

 \textbf{Answer: }D



\textbf{Solution 1}


                    \begin{figure}[H]
                        \centering
                        \includegraphics[width=0.5\textwidth]{Solution/AMC_12_12B_2021_15_Solution_Figure_1.png}
                    \end{figure}
                    

 Draw diagonals $AC$ and $AD$ to split the pentagon into three parts. We can compute the area for each triangle and sum them up at the end. For triangles $ABC$ and $ADE$, they each have area $2\cdot\frac{1}{2}\cdot\frac{4\sqrt{3}}{4}=\sqrt{3}$.* For triangle $ACD$, we can see that $AC=AD=2\sqrt{3}$ and $CD=2$. Using Pythagorean Theorem, the altitude for this triangle is $\sqrt{11}$, so the area is $\sqrt{11}$. Adding each part up, we get $2\sqrt{3}+\sqrt{11}=\sqrt{12}+\sqrt{11} \implies \boxed{\textbf{(D)} ~23}$.

 Note: Another way to find the areas would be to use the formula $A=\frac12ab\sin C$ - erringbubble

 
$\text{*}$ A more intuitive way to find this: We know $\angle B= \angle E = 120^\circ$ because of equilateral triangles, so the triangles can each be partitioned into $2$ 30-60-90 right triangles with side lengths $1,\sqrt{3},$ and $2.$ Hence, the area of the each triangle is $2\sqrt{3}\cdot 1 \cdot \frac{1}{2} = \sqrt 3.$ ~belindazhu13

 


	\item (\textit{AMC 12 12B 2021}) Let $g(x)$ be a polynomial with leading coefficient $1,$ whose three roots are the reciprocals of the three roots of $f(x)=x^3+ax^2+bx+c,$ where $1<a<b<c.$ What is $g(1)$ in terms of $a,b,$ and $c?$
$\textbf{(A) }\frac{1+a+b+c}c \qquad \textbf{(B) }1+a+b+c \qquad \textbf{(C) }\frac{1+a+b+c}{c^2}\qquad \textbf{(D) }\frac{a+b+c}{c^2} \qquad \textbf{(E) }\frac{1+a+b+c}{a+b+c}$

 \textbf{Answer: }A



\textbf{Solution 1}

Note that $f(1/x)$ has the same roots as $g(x)$, if it is multiplied by some monomial so that the leading term is $x^3$ they will be equal. We have
 \[f(1/x) = \frac{1}{x^3} + \frac{a}{x^2}+\frac{b}{x} + c\]so we can see that
 \[g(x) = \frac{x^3}{c}f(1/x)\]Therefore
 \[g(1) = \frac{1}{c}f(1) = \boxed{\textbf{(A) }\frac{1+a+b+c}c}\]




\textbf{Solution 2 (Vieta's bash)}

Let the three roots of $f(x)$ be $d$, $e$, and $f$. (Here e does NOT mean 2.7182818...)We know that $a=-(d+e+f)$, $b=de+ef+df$, and $c=-def$, and that $g(1)=1-\frac{1}{d}-\frac{1}{e}-\frac{1}{f}+\frac{1}{de}+\frac{1}{ef}+\frac{1}{df}-\frac{1}{def}$ (Vieta's). This is equal to $\frac{def-de-df-ef+d+e+f-1}{def}$, which equals $\boxed{(\textbf{A}) \frac{1+a+b+c}{c}}$. -dstanz5

 



\textbf{Solution 3 (Fakesolve)}

Because the problem doesn't specify what the coefficients of the polynomial actually are, we can just plug in any arbitrary polynomial that satisfies the constraints. Let's take $f(x) = (x+5)^3 = x^3+15x^2+75x+125$. Then $f(x)$ has a triple root of $x = -5$. Then $g(x)$ has a triple root of $-\frac{1}{5}$, and it's monic, so $g(x) = \left(x+\frac{1}{5}\right)^3 = \frac{125x^3+75x^2+15x+1}{125}$. We can see that this is $\frac{1+a+b+c}{c}$, which is answer choice $\boxed{(A)}$.

 -Darren Yao

 



\textbf{Solution 4}

If we let $p, q,$ and $r$ be the roots of $f(x)$, $f(x) = (x-p)(x-q)(x-r)$ and $g(x) = \left(x-\frac{1}{p}\right)\left(x-\frac{1}{q}\right)\left(x-\frac{1}{r}\right)$. The requested value, $g(1)$, is then
 \[\left(1-\frac{1}{p}\right)\left(1-\frac{1}{q}\right)\left(1-\frac{1}{r}\right) = \frac{(p-1)(q-1)(r-1)}{pqr}\]The numerator is $-f(1)$ (using the product form of $f(x)$ ) and the denominator is $-c$, so the answer is
 \[\frac{f(1)}{c} = \boxed{(\textbf{A}) \frac{1+a+b+c}{c}}\]

 - gting

 



\textbf{Solution 5 (Good at Guessing)}

The function $g(1) = \text{sum of coefficients}$. If it's $(x-r)(x-s)(x-t)$, then it becomes $(x-\dfrac{1}{r})(x-\dfrac{1}{s})(x-\dfrac{1}{t}).$So, $-rst$ becomes $-\dfrac{1}{rst}$, so $c$ becomes $\dfrac{1}{c}$. Also, there is a $x^3$ so the answer must include $1$. The only answer having both of these is $A$.

 -Extremelysupercooldude (Minor Latex Edits and Grammar)

 



\textbf{Solution 6}

It is well known that reversing the order of the coefficients of a polynomial turns each root into its corresponding reciprocal. Thus, a polynomial with the desired roots may be written as $cx^3+bx^2+a+1$. As the problem statement asks for a monic polynomial, our answer is 
 \[\boxed{(\textbf{A}) \frac{1+a+b+c}{c}}\]



	\item (\textit{AMC 12 12B 2021}) Let $ABCD$ be an isosceles trapezoid having parallel bases $\overline{AB}$ and $\overline{CD}$ with $AB>CD.$ Line segments from a point inside $ABCD$ to the vertices divide the trapezoid into four triangles whose areas are $2, 3, 4,$ and $5$ starting with the triangle with base $\overline{CD}$ and moving clockwise as shown in the diagram below. What is the ratio $\frac{AB}{CD}?$

                    \begin{figure}[H]
                        \centering
                        \includegraphics[width=0.5\textwidth]{Problem/AMC_12_12B_2021_17_Problem_Figure_1.png}
                    \end{figure}
                    
$\textbf{(A)}\: 3\qquad\textbf{(B)}\: 2+\sqrt{2}\qquad\textbf{(C)}\: 1+\sqrt{6}\qquad\textbf{(D)}\: 2\sqrt{3}\qquad\textbf{(E)}\: 3\sqrt{2}$

 \textbf{Answer: }B



\textbf{Solution}

Without the loss of generality, let $\mathcal T$ have vertices $A$, $B$, $C$, and $D$, with $AB = r$ and $CD = s$. Also denote by $P$ the point in the interior of $\mathcal T$.

 Let $X$ and $Y$ be the feet of the perpendiculars from $P$ to $AB$ and $CD$, respectively. Observe that $PX = \frac 8r$ and $PY = \frac 4s$. Now using the formula for the area of a trapezoid yields
 \[14 = \frac12\cdot XY\cdot (AB+CD) = \frac12\left(\frac 8r + \frac 4s\right)(r+s) = 6 + 2\cdot\frac rs + 4\cdot\frac sr.\]Thus, the ratio $\rho := \frac rs$ satisfies $\rho + 2\rho^{-1} = 4$; solving yields $\rho = \boxed{\textbf{(B)}\: 2+\sqrt{2}}$.

 (Observe that the given areas of $3$ and $5$ are irrelevant to the ratio $\frac{AB}{CD}$.)

 



\textbf{Solution 2}

Let $b_1$ be the bottom base, $b_2$ be the top base, $h_1$ be the height of the bottom triangle, $h_2$ be the height of the top triangle.Thus, $b_1h_1 = 8, b_2h_2 = 4, (b_1+b_2)(h_1+h_2) = 28,$ so $b_1h_2 + b_2h_1 = 16.$ Let $b_2 = 1, h_2 = 4,$ so we get $b_1h_1 = 8, 4b_1+h_1 = 16.$This gives us a quadratic in $b_1,$ ie. $4b_1^2+8=16b_1,$ so $b_1 = \boxed{\textbf{(B)}\: 2+\sqrt{2}}.$

 - Solution by MathAwesome123, added by ccx09

 


 


	\item (\textit{AMC 12 12B 2021}) Let $z$ be a complex number satisfying $12|z|^2=2|z+2|^2+|z^2+1|^2+31.$ What is the value of $z+\frac 6z?$
$\textbf{(A) }-2 \qquad \textbf{(B) }-1 \qquad \textbf{(C) }\frac12\qquad \textbf{(D) }1 \qquad \textbf{(E) }4$

 \textbf{Answer: }A



\textbf{Solution 1}

Using the fact $z\bar{z}=|z|^2$, the equation rewrites itself as
 \begin{align*} 12z\bar{z}&=2(z+2)(\bar{z}+2)+(z^2+1)(\bar{z}^2+1)+31 \\ -12z\bar{z}+2z\bar{z}+4(z+\bar{z})+8+z^2\bar{z}^2+(z^2+\bar{z}^2)+32&=0 \\ \left((z^2+2z\bar{z}+\bar{z}^2)+4(z+\bar{z})+4\right)+\left(z^2\bar{z}^2-12z\bar{z}+36\right)&=0 \\ (z+\bar{z}+2)^2+(z\bar{z}-6)^2&=0. \end{align*}
As the two quantities in the parentheses are real, both quantities must equal $0$ so 
 \[z+\frac6z=z+\bar{z}=\boxed{\textbf{(A) }-2}.\]




\textbf{Solution 2}

Let $z = a + bi$, $z^2 = a^2-b^2+2abi$

 By the equation given in the problem

 
 \[12(a^2+b^2) = 2((a+2)^2 + b^2) + ((a^2-b^2+1)^2 + (2ab)^2) + 31\]

 \[12a^2 + 12b^2 = 2a^2 + 8a + 8 + 2b^2 + a^4 + b^4 + 1 + 2a^2 - 2b^2 - 2a^2b^2 + 4a^2b^2 + 31\]

 \[a^4 + b^4 - 8a^2 - 12b^2 + 2a^2b^2 + 8a + 40 = 0\]

 \[(a^2+b^2)^2 - 12(a^2+b^2) + 4(a^2 + 2a + 1) + 36=0\]

 \[(a^2 + b^2 - 6)^2 + 4(a+1)^2 = 0\]

 Therefore, $a^2 + b^2 - 6 = 0$ and $a+1 = 0$
$a = -1$, $b^2 = 6-1 = 5$, $b = \sqrt{5}$

 \[z + \frac{6}{z} = \frac{ a^2 - b^2 + 6 + 2abi }{ a+bi } = \frac{ 1 - 5 + 6 + 2(-1)\sqrt{5} i }{ -1 + i \sqrt{5} } = \frac{ 2 - 2i \sqrt{5} }{-1 + i \sqrt{5}} = \boxed{\textbf{(A)} -2}\]

 should be $b =  \sqrt{5}$ or $b =  -\sqrt{5}$ but it won't affect the result. ~eddybai

 



\textbf{Solution 3}

Let $x = z + \frac{6}{z}$. Then $z = \frac{x \pm \sqrt{x^2-24}}{2}$. From the answer choices， we know that $x$ is real and $x^2<24$, so $z = \frac{x \pm i\sqrt{24-x^2}}{2}$. Then we have
 \[|z|^2 = 6\]
 \[|z+2|^2 = (\frac{x}{2} + 2)^2 + \frac{24-x^2}{4} = 2x+10\]
 \[|z^2+1|^2 = |xz -6 +1|^2 = \left(\frac{x^2}{2}-5\right)^2 + \frac{x^2(24-x^2)}{4} = x^2 +25\]Plugging the above back to the original equation, we have
 \[12*6 = 2(2x+10) + x^2 + 25 + 31\]
 \[(x+2)^2 = 0\]So $x = \boxed{\textbf{(A) }-2}$. 

 



\textbf{Solution 4 (Funny Observations)}

There are actually several ways to see that $|z|^2 = 6.$ I present two troll ways of seeing it, and a legitimate way of checking.

 Rewrite using $w \overline{w} = |w|^2$
$12z \overline{z} + 2(z+2)(\overline{z} + 2) + (z^2+1)(\overline{z}^2+1)+31$$12 z \overline{z} = 2z \overline{z} + 4z + 4 \overline{z} + 8 + z^2 \overline{z}^2+z^2+\overline{z}^2 + 1 + 31.$$12 z \overline{z} = 4(z + \overline{z}) + (z \overline{z})^2 + (z + \overline{z})^2 + 40.$

 Symmetric in $z$ and $\overline{z},$ so if $w$ is a sol, then so is $\overline{w}$

 TROLL OBSERVATION \#1: ALL THE ANSWERS ARE REAL. THUS, $z + \frac{6}{z} \in \mathbb{R},$ which means they must be conjugates and so $|z|^2 = 6.$

 TROLL OBSERVATION \#2: Note that $z+\frac{6}{z} = \overline{z} + \frac{6}{\overline{z}}$ because either solution must give the same answer! which means that $|z|^2 = 6.$ 

 Alternatively, you can check:Let $a = w + \overline{w} \in \mathbb{R},$ and $r = |w|^2 \in \mathbb{R}.$ Thus, we have $a^2+4a+40+r^2-12r=0,$ and the discriminant of this must be nonnegative as $a$ is real. Thus, $16-4(40+r^2-12r) \geq 0$ or $(r-6)^2 \leq 0,$ which forces $r = 6,$ as claimed.

 Thus, we plug in $z \overline{z} = 6,$ and get: $72 = 4(z + \overline{z}) + 76 + (z + \overline{z})^2,$ ie. $(z+\overline{z})^2 + 4(z + \overline{z}) + 4 = 0,$ or $(z+\overline{z} + 2)^2 = 0,$ which means $z + \overline{z} = \boxed{\textbf{(A) }-2}$ and that's our answer since we know $\overline{z} = 6 / z$

 - ccx09

 



\textbf{Solution 5}

Observe that all the answer choices are real. Therefore, $z$ and $\frac{6}{z}$ must be complex conjugates as this is the only way for both their sum (one of the answer choices) and their product ($6$) to be real. Thus $|z|=|\frac{6}{z}|=\sqrt{6}$. We will test all the answer choices, starting with $\textbf{(A)}$. Suppose the answer is $\textbf{(A)}$. If $z+\frac{6}{z}=-2$ then $z^{2}+2z+6=0$ and $z=\frac{-2\pm\sqrt{4-24}}{2}=-1\pm\sqrt{5}i$. Note that if $z=-1+\sqrt{5}i$ works, then so does $-1-\sqrt{5}i$. It is relatively easy to see that if $z=-1+\sqrt{5}i$, then $12|z|^{2}=12\cdot 6=72, 2|z+2|^{2}=2|1+\sqrt{5}i|=2\cdot 6=12, |z^{2}+1|^{2}=|-3-2\sqrt{5}i|^{2}=29,$ and $72=12+29+31$. Thus the condition 
 \[12|z|^{2}=2|z+2|^{2}+|z^{2}+1|^{2}+31\] is satisfied for $z+\frac{6}{z}=-2$, and the answer is $\boxed{\textbf{(A) }-2}$.

 



\textbf{Solution 6}

Using $z\bar{z}=|z|^2$, we have\begin{align*}12z\bar{z} = 2(z + 2)(\bar{z} + 2) + (z^2 + 1)(\bar{z}^2 + 1) + 31 &\implies z^2\bar{z}^2 + z^2 + \bar{z}^2 + 1 + 2z\bar{z} + 4z + 4\bar{z} + 8 + 31 - 12z\bar{z} = 0 \\&\implies  z^2\bar{z}^2 - 10z\bar{z} + (z^2 + \bar{z}^2) + 4(z + \bar{z}) + 40 = 0\end{align*}
Let $p = z\bar{z}$ and $s = z + \bar{z}$. Then we get\begin{align*}p^2 - 10p + s^2 - 2p + 4s + 40 = 0 \implies p^2 - 12p + s^2 + 4s + 40 = 0\end{align*}
Completing the square, we get
 \[(p-6)^2 - 36 + (s+2)^2 - 4 + 40 = (p-6)^2 + (s+2)^2 = 0\]Therefore, $p = 6$ and $s = -2$. So, $z + \bar{z} = -2$ and $z\bar{z} = 6$. Plugging into $z + \frac{6}{z}$, we get
 \[z + \frac{6}{z} = z + \frac{6\bar{z}}{z\bar{z}} = z + \frac{6\bar{z}}{6} = z + \bar{z} = \boxed{\textbf{(A) }-2}\]



	\item (\textit{AMC 12 12B 2021}) Two fair dice, each with at least $6$ faces are rolled. On each face of each die is printed a distinct integer from $1$ to the number of faces on that die, inclusive. The probability of rolling a sum of $7$ is $\frac34$ of the probability of rolling a sum of $10,$ and the probability of rolling a sum of $12$ is $\frac{1}{12}$. What is the least possible number of faces on the two dice combined?


 $\textbf{(A) }16 \qquad \textbf{(B) }17 \qquad \textbf{(C) }18\qquad \textbf{(D) }19 \qquad \textbf{(E) }20$

 \textbf{Answer: }B



\textbf{Solution 1}

Suppose the dice have $a$ and $b$ faces, and WLOG $a\geq{b}$. Since each die has at least $6$ faces, there will always be $6$ ways to sum to $7$. As a result, there must be $\frac{4}{3}\cdot6=8$ ways to sum to $10$. There are at most nine distinct ways to get a sum of $10$, which are possible whenever $a,b\geq{9}$. To achieve exactly eight ways, $b$ must have $8$ faces, and $a\geq9$. Let $n$ be the number of ways to obtain a sum of $12$, then $\frac{n}{8a}=\frac{1}{12}\implies n=\frac{2}{3}a$. Since $b=8$, $n\leq8\implies a\leq{12}$. In addition to $3\mid{a}$, we only have to test $a=9,12$, of which both work. Taking the smaller one, our answer becomes $a+b=9+8=\boxed{\textbf{(B)}\ 17}$.

 



\textbf{Solution 2}

Suppose the dice have $a$ and $b$ faces, and WLOG $a\geq{b}$. Note that if $a+b=12$ since they are both $6$, there is one way to make $12$, and incrementing $a$ or $b$ by one will add another way. This gives us the probability of making a 12 as 
 \[\frac{a+b-11}{ab}=\frac{1}{12}\]Cross-multiplying, we get
 \[12a+12b-132=ab\]Simon's Favorite Factoring Trick now gives
 \[(a-12)(b-12)=12\]This narrows the possibilities down to 3 ordered pairs of $(a,b)$, which are $(13,24)$, $(6,10)$, and $(8,9)$.  We can obviously ignore the first pair and test the next two straightforwardly.  The last pair yields the answer:
 \[\frac{6}{72}=\frac{3}{4}\left(\frac{9+8-9}{72}\right)\]The answer is then $a+b=8+9=\boxed{\textbf{(B)}\ 17}$.

 



\textbf{Solution 2.1}

notice that if a or b >= 12 then this relationship doesn't stand.the value of a and b should be discussed.Hence the final solution should be $(6,10)$, and $(8,9)$ and $(12,b)$ (b<=11)
 \[\frac{min(a,11)+min(b,11)-11}{ab}=\frac{1}{12}\]and to check P(7) and P(10), simply use 
 \[\frac{min(a,6)+min(b,6)-6}{ab}\] and 
 \[\frac{min(a,9)+min(b,9)-9}{ab}\]~eddybai

 



\textbf{Solution 3 (Logic)}

Notice that
 \begin{align*}     7&=1+6 \\     &\ \ \ \ \ \ \ \ \ \vdots \\     &=6+1 \end{align*}
are the only cases that could possibly form 7. In another words, regardless of the number of faces for each dice, 6 is the number of cases that could form a rolling sum of 7.

 
Because the value of the total combination for both are the same, we could infer that the number of cases of having a rolling sum of 10 is $\frac{4}{3}\cdot6=8$. With the number 8, we could also deduce that one dice has 8 sides and the other has at least 9 sides. Thence trial and error could be utilized.

 
Since 9 is the next smallest number, the case could be tested.
 \begin{align*}     12&=3+9 \\     &=4+8 \\     &\ \ \ \ \ \ \ \ \ \vdots \\     &=8+4 \end{align*}
$\frac{6}{8\cdot9}=\frac{1}{12}$ is also true. Therefore, $8+9=\boxed{\textbf{(B) }17}$




\textbf{Solution 4}

Let the two dice have at least six faces each, and let their face counts sum to $N$. We seek the minimum possible $N$ such that $p(\text{sum}=12)=\frac{1}{12}$ and $p(\text{sum}=7)=\frac{3}{4}p(\text{sum}=10)$.

 First try $N=16$. Then the possible face-count pairs $(m,n)$ with $m \leq n$ are $(6,10)$, $(7,9)$, $(8,8)$. For $p(\text{sum}=12)=\frac{1}{12}$ to be possible, the total outcomes $mn$ must be a multiple of 12. But $7\cdot 9=63$ and $8\cdot 8=64$ are not multiples of 12, so $(7,9)$ and $(8,8)$ are impossible.  

 For $(6,10)$, there are 60 outcomes. The combinations that sum to 12 are (2,10), (3,9), (4,8), (5,7), (6,6), five in total, giving $p(12)=\frac{5}{60}=\frac{1}{12}$.  The combinations that sum to 10 are (1,9), (2,8), (3,7), (4,6), (5,5), (6,4), six in total, giving $p(10)=\frac{6}{60}=\frac{1}{10}$.  The combinations that sum to 7 are (1,6), (2,5), (3,4), (4,3), (5,2), (6,1), also six in total, giving $p(7)=\frac{6}{60}=\frac{1}{10}$.  So $p(7)=p(10)$, which does not satisfy the condition. Hence $N=16$ is impossible.  

 Next try $N=17$. The possible pairs are $(6,11)$, $(7,10)$, $(8,9)$. Again, $mn$ must be a multiple of 12, but $6\cdot 11=66$ and $7\cdot 10=70$ are not; discard these.  

 For $(8,9)$, there are 72 outcomes.  The combinations that sum to 12 are (3,9), (4,8), (5,7), (6,6), (7,5), (8,4), six in total, so $p(12)=\frac{6}{72}=\frac{1}{12}$.  The combinations that sum to 10 are (1,9), (2,8), (3,7), (4,6), (5,5), (6,4), (7,3), (8,2), eight in total, so $p(10)=\frac{8}{72}=\frac{1}{9}$.  The combinations that sum to 7 are (1,6), (2,5), (3,4), (4,3), (5,2), (6,1), six in total, so $p(7)=\frac{6}{72}=\frac{1}{12}$.  

 Thus $p(12)=\frac{1}{12}$ and $p(7)=\frac{1}{12}=\frac{3}{4}\cdot\frac{1}{9}=\frac{3}{4}p(10)$.  

 Therefore $N=17$ works, and since $N=16$ is impossible, the minimum is $\boxed{(\textbf{B})\ 17}$.

 


	\item (\textit{AMC 12 12B 2021}) Let $Q(z)$ and $R(z)$ be the unique polynomials such that
 \[z^{2021}+1=(z^2+z+1)Q(z)+R(z)\]and the degree of $R$ is less than $2.$ What is $R(z)?$
$\textbf{(A) }{-}z \qquad \textbf{(B) }{-}1 \qquad \textbf{(C) }2021\qquad \textbf{(D) }z+1 \qquad \textbf{(E) }2z+1$

 \textbf{Answer: }A



\textbf{Solution 1 (Difference of Cubes)}

Let $z=s$ be a root of $z^2+z+1$ so that $s^2+s+1=0.$ It follows that 
 \[(s-1)\left(s^2+s+1\right)=s^3-1=0,\] from which $s^3=1,$ but $s\neq1.$

 Note that
 \begin{align*} s^{2021}+1 &= s^{3\cdot673+2}+1 \\ &= (s^3)^{673}\cdot s^2+1 \\ &= s^2+1 \\ &= \left(s^2+s+1\right)-s \\ &= -s. \end{align*}
Since $z^{2021}+1=-z$ for each root $z=s$ of $z^2+z+1,$ the remainder when $z^{2021}+1$ is divided by $z^2+z+1$ is $R(z)=\boxed{\textbf{(A) }{-}z}.$




\textbf{Solution 2 (Finds Q(z) Using Patterns)}

Note that the equation above is in the form of polynomial division, with $z^{2021}+1$ being the dividend, $z^2+z+1$ being the divisor, and $Q(x)$ and $R(x)$ being the quotient and remainder respectively. Since the degree of the dividend is $2021$ and the degree of the divisor is $2$, that means the degree of the quotient is $2021-2 = 2019$. Note that $R(x)$ can't influence the degree of the right hand side of this equation since its degree is either $1$ or $0$. Since the coefficients of the leading term in the dividend and the divisor are both $1$, that means the coefficient of the leading term of the quotient is also $1$. Thus, the leading term of the quotient is $z^{2019}$. Multiplying $z^{2019}$ by the divisor gives $z^{2021}+z^{2020}+z^{2019}$. We have our $z^{2021}$ term but we have these unnecessary terms like $z^{2020}$. We can get rid of these terms by adding $-z^{2018}$ to the quotient to cancel out these terms, but this then gives us $z^{2021}-z^{2018}$. Our first instinct will probably be to add $z^{2017}$, but we can't do this as although this will eliminate the $-z^{2018}$ term, it will produce a $z^{2019}$ term. Since no other term of the form $z^n$ where $n$ is an integer less than $2017$ will produce a $z^{2019}$ term when multiplied by the divisor, we can't add $z^{2017}$ to the quotient. Instead, we can add $z^{2016}$ to the coefficient to get rid of the $-z^{2018}$ term. Continuing this pattern, we get the quotient as 
 \[z^{2019}-z^{2018}+z^{2016}-z^{2015}+\cdots-z^2+1.\]The last term when multiplied with the divisor gives $z^2+z+1$. This will get rid of the $-z^2$ term but will produce the expression $z+1$, giving us the dividend as $z^{2021}+z+1$. Note that the dividend we want is of the form $z^{2021}+1$. Therefore, our remainder will have to be $-z$ in order to get rid of the $z$ term in the expression and give us $z^{2021}+1$, which is what we want. Therefore, the remainder is $\boxed{\textbf{(A) }{-}z}.$




\textbf{Solution 3 (Modular Arithmetic in Polynomials)}

Note that
 \[z^3-1\equiv 0\pmod{z^2+z+1}\]so if $F(z)$ is the remainder when dividing by $z^3-1$,
 \[F(z)\equiv R(z)\pmod{z^2+z+1}.\]Now,
 \[z^{2021}+1= (z^3-1)(z^{2018} + z^{2015} + \cdots + z^2) + z^2+1\]So $F(z) = z^2+1$, and
 \[R(z)\equiv F(z) \equiv -z\pmod{z^2+z+1}\]The answer is $\boxed{\textbf{(A) }{-}z}.$




\textbf{Solution 4 (Complex Numbers)}

One thing to note is that $R(z)$ takes the form of $Az + B$ for some constants $A$ and $B.$Note that the roots of $z^2 + z + 1$ are part of the solutions of $z^3 -1 = 0.$They can be easily solved with roots of unity:
 \begin{align*} z^3 &= 1 \\ z^3 &= e^{i 0} \\ z &= e^{i 0}, e^{i \frac{2\pi}{3}}, e^{i\left(-\frac{2\pi}{3}\right)}. \end{align*}
Obviously the right two solutions are the roots of $z^2 + z + 1 = 0.$We substitute $e^{i \frac{2\pi}{3}}$ into the original equation, and $z^2 + z + 1$ becomes $0.$ Using De Moivre's Theorem, we get
 \begin{align*} e^{i\frac{4042\pi}{3}} + 1 &= A \cdot e^{i \frac{2\pi}{3}} + B \\ e^{i\frac{4\pi}{3}} + 1 &= A \cdot e^{i \frac{2\pi}{3}} + B. \end{align*}
Expanding into rectangular complex number form:
 \[\frac{1}{2} - \frac{\sqrt{3}}{2} i = \left(-\frac{1}{2}A + B\right) + \frac{\sqrt{3}}{2} A i.\]Comparing the real and imaginary parts, we get $(A,B) = (-1,0).$ So, the answer is $\boxed{\textbf{(A) }{-}z}.$




\textbf{Solution 5 (Complex Numbers but Easier)}

We have motivation to get rid of the $Q(z)$ term by subtituting in either $e^{\frac{2\pi i}{3}}$ or $e^{\frac{4\pi i}{3}}$ for $z$, as this sets $z^2+z+1$ equal to $0$. Doing so, 
 \begin{align*} \left(e^{\frac{2\pi i}{3}}\right)^{2021} + 1 &= R\left(e^{\frac{2\pi i}{3}}\right) \\ e^{\frac{4042\pi i}{3}} + 1 &= R\left(e^{\frac{2\pi i}{3}}\right) \\ e^{\frac{4\pi i}{3}} + 1 &= R\left(e^{\frac{2\pi i}{3}}\right) \\ \frac{1}{2}-\frac{\sqrt{3}}{2}i &= R\left(-\frac{1}{2} + \frac{\sqrt{3}}{2} i \right). \end{align*}
This immediately eliminates choices $\textbf{(B)}$ and $\textbf{(C)},$ as they are constants. Upon a quick check, $R(z) = 2z+1$ and $z + 1$ both don't work, as the sign of $\frac{\sqrt{3}}{2}$ has to be negative. Checking $R(z)$ (or realizing that it's the only one that actually works), we see that $\boxed{\textbf{(A) }{-}z}$ is the right answer. 

 -Benedict T (countmath1)

 



\textbf{Solution 6 (More Complex Numbers)}

Since $z^{2021}+1=(z^2+z+1)Q(z)+R(z),$ we can plug in either $\text{cis} (240^{\circ})$ or $\text{cis} (240^{\circ})$. Assuming $R(z)=ax+b,$ we have $a(\text{cis}(240^{\circ}))+b=\frac{1}{2}+\frac{i\sqrt{3}}{2}$ and $a(\text{cis}(120^{\circ}))+b=\frac{1}{2}-\frac{i\sqrt{3}}{2}.$ Solving gives us $a=-1$ and $b=0$, thus $R(z)=\boxed{\textbf{(A) }{-}z}.$




\textbf{Solution 7 (Factorization)}

Notice that $(z^{2018} + z^{2015} + z^{2012} + \cdots + z^5 + z^2)(z^2 + z + 1) = z^{2020} + z^{2019} + \cdots + z^{4} + z^3 + z^2$, so 
 \begin{align*} z^{2021} + 1 = z^{2021} - 1 + 2 &= (z-1)( z^{2020} + z^{2019} + \cdots + z^{2} + z + 1 ) + 2 \\ &= (z-1)[(z^{2018} + z^{2015} + z^{2012} + \cdots + z^5 + z^2)(z^2 + z + 1) + z + 1] + 2 \\ &= (z-1)(z^{2018} + z^{2015} + z^{2012} + \cdots + z^5 + z^2)(z^2 + z + 1) + (z-1)(z+1) + 2 \\ &= (z-1)(z^{2018} + z^{2015} + z^{2012} + \cdots + z^5 + z^2)(z^2 + z + 1) + z^2 + 1 \\ &= (z-1)(z^{2018} + z^{2015} + z^{2012} + \cdots + z^5 + z^2)(z^2 + z + 1) + z^2 + z + 1 - z \\ &= (z^2 + z + 1)[(z^{2018} + z^{2015} + z^{2012} + \cdots + z^5 + z^2)(z-1) + 1] - z. \end{align*}
Therefore, we have 
 \[Q(z) = (z^{2018} + z^{2015} + z^{2012} + \cdots + z^5 + z^2)(z-1) + 1,\] from which $R(z) = \boxed{\textbf{(A) }{-}z}$.

 



\textbf{Solution 8}

Let $P(z)=z^{2021}+1$ and $r$ be a root of $z^2+z+1$, then note that

 
 \begin{align*} &r^{3k} = 1 \\ &r^{3k+1} = r \\ &r^{3k+2} = -r-1 \\ \end{align*}


 for all non-negative integers $k$. Since
 \[P(r)=(r^2+r+1)Q(r)+R(r)=R(r)\]and $P(r)=r^{2021}+1=(-r - 1)+1=-r$, $R(r)=-r$. Therefore, $\boxed{\textbf{(A) }{-}z}$ is the only feasible option.

 



\textbf{Solution 9 (Modular Arithmetic in less than a minute, basically a slightly modified solution 1)}

Notice that $z^{3} - 1 \equiv 0 \pmod{z^2 + z + 1}$ because $z^3 - 1 = (z-1)(z^{2} + z + 1).$ So we can add this to $z^{2021} + 1$: 
 \[z^{2021} + z^{3} \equiv z^{2021} + 1 \pmod{z^2 + z + 1}\]
 \[\implies z^{3}(z^{2018} + 1) \equiv 1*(z^{2018} + 1) \pmod{z^2 + z + 1}\] We can generalize this by saying $z^{2021 - 3n} + 1 \equiv z^{2021} + 1 \pmod{z^2 + z + 1}$ as long as $2021 - 3n$ is positive. So we can minimize $2021 - 3n$ by setting $n = 673.$ Then we get: 
 \[z^{2} + 1 \equiv z^{2021} + 1 \pmod{z^2 + z + 1}\] And this is clearly $\boxed{\textbf{(A) }{-}z}$. 

 



\textbf{Solution 10 (Easy Modular arithmetic, Similar to solution 9)}

Similarly $z^{3} - 1 \equiv 0 \pmod{z^2 + z + 1}$ because $z^3 - 1 = (z-1)(z^{2} + z + 1).$

 So,

 
 \[z^3\equiv1\pmod{z^2 + z + 1}\]

 \[\implies R(z) \equiv z^{2021} +1\equiv (z^{3})^{673} \cdot z^{2}+1 \equiv z^{2} +1  \pmod{z^2 + z + 1}\]

 
Meaning,

 
 \[R(z)\equiv z^{2} +1 \pmod{z^2 + z + 1}\]

 Notice $z^{2}\equiv -(z+1)  \pmod{z^2 + z + 1}\implies z^2+1\equiv -(z+1)+1\equiv -z  \pmod{z^2 + z + 1}$ so

 
 \[R(z)\equiv -z  \pmod{z^2 + z + 1}\]

 Gives our answer  $\boxed{\textbf{(A) }{-}z}$. 

 


 


	\item (\textit{AMC 12 12B 2021}) Let $S$ be the sum of all positive real numbers $x$ for which
 \[x^{2^{\sqrt2}}=\sqrt2^{2^x}.\]Which of the following statements is true?


 $\textbf{(A) }S<\sqrt2 \qquad \textbf{(B) }S=\sqrt2 \qquad \textbf{(C) }\sqrt2<S<2\qquad \textbf{(D) }2\le S<6 \qquad \textbf{(E) }S\ge 6$

 \textbf{Answer: }D



\textbf{Solution 1}

Note that
 \begin{align*} x^{2^{\sqrt{2}}} &= {\sqrt{2}}^{2^x} \\ 2^{\sqrt{2}} \log_2 x &= 2^{x} \log_2 \sqrt{2}. \end{align*}
(At this point we see by inspection that $x=\sqrt{2}$ is a solution.)

 We simplify the RHS, then take the base-$2$ logarithm for both sides:
 \begin{align*} 2^{\sqrt{2}} \log_2 x &= 2^{x-1} \\ \log_2{\left(2^{\sqrt{2}} \log_2 x\right)} &= x-1 \\ \sqrt{2} + \log_2 \log_2 x &= x-1 \\ \log_2 \log_2 x &= x - 1 - \sqrt{2}. \end{align*}
The RHS is a line; the LHS is a concave curve that looks like a logarithm and has $x$ intercept at $(2,0).$

 There are at most two solutions, one of which is $\sqrt{2}.$ But note that at $x=2,$ we have $\log_2 \log_2 {2} = 0 > 2 - 1 - \sqrt{2},$ meaning that the log log curve is above the line, so it must intersect the line again at a point $x > 2.$ Now we check $x=4$ and see that $\log_2 \log_2 {4} = 1 < 4 - 1 - \sqrt{2},$ which means at $x=4$ the line is already above the log log curve. Thus, the second solution lies in the interval $(2,4).$The answer is $\boxed{\textbf{(D) }2\le S<6}.$




\textbf{Solution 2}

We rewrite the right side without using square roots, then take the base-$2$ logarithm for both sides:
 \begin{align*} x^{2^{\sqrt2}}&=\left(2^\frac12\right)^{2^x} \\ x^{2^{\sqrt2}}&=2^{\frac12\cdot2^x} \\ x^{2^{\sqrt2}}&=2^{2^{x-1}} \\ \log_2{\left(x^{2^{\sqrt2}}\right)}&=\log_2{\left(2^{2^{x-1}}\right)} \\ 2^{\sqrt2}\log_2{x}&=2^{x-1}. \hspace{20mm} (*) \end{align*}
By observations, $x=\sqrt2$ is one solution. Graphing $f(x)=2^{\sqrt2}\log_2{x}$ and $g(x)=2^{x-1},$ we conclude that $(*)$ has two solutions, with the smaller solution $x=\sqrt2.$ We construct the following table of values: 
 \[\begin{array}{c|c|c|c} & & & \\ [-2ex] \boldsymbol{x} & \boldsymbol{f(x)=2^{\sqrt2}\log_2{x}} & \boldsymbol{g(x)=2^{x-1}} & \textbf{Comparison} \\ [1ex] \hline & & & \\ [-1ex] 1 & 0 & 1 & \\ [1ex] \sqrt2 & 2^{\sqrt2-1} & 2^{\sqrt2-1} & f\left(\sqrt2\right)=g\left(\sqrt2\right) \\ [1ex] 2 & 2^{\sqrt2} & 2 & f(2)>g(2) \\ [1ex] 3 & 2^{\sqrt2}\log_2{3} & 2^2 & \\ [1ex] 4 & 2^{\sqrt2+1} & 2^3 & f(4)<g(4) \\ [1ex] \end{array}\]Let $x=t$ be the larger solution. Since exponential functions outgrow logarithmic functions, we have $f(x)<g(x)$ for all $x>t.$ By the \textbf{Intermediate Value Theorem}, we get $t\in(2,4),$ from which 
 \[S=\sqrt2+t\in\left(\sqrt2+2,\sqrt2+4\right).\]Finally, approximating with $\sqrt2\approx1.414$ results in $\boxed{\textbf{(D) }2\le S<6}.$

 The graphs of $y=f(x)$ and $y=g(x)$ are shown below:
                    \begin{figure}[H]
                        \centering
                        \includegraphics[width=0.5\textwidth]{Solution/AMC_12_12B_2021_21_Solution_Figure_1.png}
                    \end{figure}
                    ~MRENTHUSIASM

 



\textbf{Solution 3 (No Logarithm)}

While utilizing $\log$ may seem conventional, graphing may also be used. Notice that $x^{2^{\sqrt2}}$ is a U-shaped, differentiable curve. Moreover, $\sqrt2^{2^x}$ is an exponentially increasing function. Furthermore, we could notice that the two graphs meet at $x=\sqrt{2}$.
 \[\begin{array}{c|c|c|c}     x & x^{2^{\sqrt2}} & \sqrt2^{2^x} & \text{Comparison} \\     \hline     0 & 0 & \sqrt{2} & x^{2^{\sqrt2}} < \sqrt2^{2^x} \\     1 & 2<2^{\sqrt{2}}<4 & 2 & x^{2^{\sqrt2}} < \sqrt2^{2^x} \\     \sqrt{2} & {\sqrt{2}}^{2^{\sqrt2}} & {\sqrt{2}}^{2^{\sqrt2}} & x^{2^{\sqrt2}} = \sqrt2^{2^x} \\     2 & 4<2^{2^{\sqrt2}}<16 & 4 & x^{2^{\sqrt2}} > \sqrt2^{2^x} \\     3 & 9<3^{2^{\sqrt2}}<81 & 16 & NA \\     4 & 16<4^{2^{\sqrt2}}<256 & 256 & x^{2^{\sqrt2}} < \sqrt2^{2^x} \\ \end{array}\]Because the $x^2<x^{2^{\sqrt2}}<x^4$ and $\sqrt2^{2^x}$ is an exponentially increasing function, $x^{2^{\sqrt2}}$ can never meet or catch up $\sqrt2^{2^x}$ from $x\ge4$. In another words, only two intersections, or the roots of the equation, exists. The roots are $\sqrt{2}$ and a constant in the interval $2<r_2<4$. Let $S=\sqrt{2}+r_2$. It is evident that the possible range for $S$ is $\boxed{\textbf{(D) }2\le S<6}$.

 


	\item (\textit{AMC 12 12B 2021}) Arjun and Beth play a game in which they take turns removing one brick or two adjacent bricks from one "wall" among a set of several walls of bricks, with gaps possibly creating new walls. The walls are one brick tall. For example, a set of walls of sizes $4$ and $2$ can be changed into any of the following by one move: $(3,2),(2,1,2),(4),(4,1),(2,2),$ or $(1,1,2).$

                    \begin{figure}[H]
                        \centering
                        \includegraphics[width=0.5\textwidth]{Problem/AMC_12_12B_2021_22_Problem_Figure_1.png}
                    \end{figure}
                    


 Arjun plays first, and the player who removes the last brick wins. For which starting configuration is there a strategy that guarantees a win for Beth?


 $\textbf{(A) }(6,1,1) \qquad \textbf{(B) }(6,2,1) \qquad \textbf{(C) }(6,2,2)\qquad \textbf{(D) }(6,3,1) \qquad \textbf{(E) }(6,3,2)$

 \textbf{Answer: }B



\textbf{Solution 1}

We say that a game state is an N-position if it is winning for the next player (the player to move), and a P-position if it is winning for the other player. We are trying to find which of the given states is a P-position.

 First we note that symmetrical positions are P-positions, as the second player can win by mirroring the first player's moves. It follows that $(6, 1, 1)$ is an N-position, since we can win by moving to $(2, 2, 1, 1)$; this rules out $\textbf{(A)}$. We next look at $(6, 2, 1)$. The possible next states are
 \[(6, 2), \enskip (6, 1, 1), \enskip (6, 1), \enskip (5, 2, 1), \enskip (4, 2, 1, 1), \enskip (4, 2, 1), \enskip (3, 2, 2, 1), \enskip (3, 2, 1, 1), \enskip (2, 2, 2, 1).\]None of these are symmetrical, so we might reasonably suspect that they are all N-positions. Indeed, it just so happens that for all of these states except $(6, 2)$ and $(6, 1)$, we can win by moving to $(2, 2, 1, 1)$; it remains to check that $(6, 2)$ and $(6, 1)$ are N-positions.

 To save ourselves work, it would be nice if we could find a single P-position directly reachable from both $(6, 2)$ and $(6, 1)$. We notice that $(3, 2, 1)$ is directly reachable from both states, so it would suffice to show that $(3, 2, 1)$ is a P-position. Indeed, the possible next states are
 \[(3, 2), \enskip (3, 1, 1), \enskip (3, 1), \enskip (2, 2, 1), \enskip (2, 1, 1, 1), \enskip (2, 1, 1),\]which allow for the following refutations:

 
 \begin{align*} &(3, 2) \to (2, 2), && &&(3, 1, 1) \to (1, 1, 1, 1), && &&(3, 1) \to (1, 1), \\ &(2, 2, 1) \to (2, 2), && &&(2, 1, 1, 1) \to (1, 1, 1, 1), && &&(2, 1, 1) \to (1, 1). \end{align*}


 Hence, $(3, 2, 1)$ is a P-position, so $(6, 2)$ and $(6, 1)$ are both N-positions, along with all other possible next states from $(6, 2, 1)$ as noted before. Thus, $(6, 2, 1)$ is a P-position, so our answer is $\boxed{\textbf{(B)}}$. (For completeness, we could also rule out $\textbf{(C)}$, $\textbf{(D)}$ and $\textbf{(E)}$ as in Solution 2.)

 -Note: In general, this game is very complicated. For example $(8, 7, 5, 3, 2)$ is winning for the first player but good luck showing that.

 



\textbf{Solution 2 (Elimination)}

$(6,1,1)$ can be turned into $(2,2,1,1)$ by Arjun, which is symmetric, so Beth will lose.

 $(6,3,1)$ can be turned into $(3,1,3,1)$ by Arjun, which is symmetric, so Beth will lose.

 $(6,2,2)$ can be turned into $(2,2,2,2)$ by Arjun, which is symmetric, so Beth will lose.

 $(6,3,2)$ can be turned into $(3,2,3,2)$ by Arjun, which is symmetric, so Beth will lose.

 That leaves $(6,2,1)$ or $\boxed{\textbf{(B)}}$.

 



\textbf{Solution 3 (Nim-values)}

Let the nim-value of the ending game state, where someone has just removed the final brick, be $0$. Then, any game state with a nim-value of $0$ is losing. It is well-known that the nim-value of a supergame (a combination of two or more individual games) is the binary xor function on the nim-values of the individual games that compose the supergame. Therefore, we calculate the nim-values of the states with a single wall up to $6$ bricks long (since the answer choices only go up to $6$). 

 First, the game with $1$ brick has a nim-value of $1$. 

 Similarly, the game with $2$ bricks has a nim-value of $2$.

 Next, we consider a $3$ brick wall. After the next move, the possible resulting game states are $1$ brick, a $2$ brick wall, or $2$ separate bricks. The first two options have nim-values of $1$ and $2$. The final option has a nim-value of $1\oplus 1 = 0$, so the nim-value of this game state is $3$. 

 Next, the $4$ brick wall. The possible states are a $2$ brick wall, a $3$ brick wall, a $2$ brick wall and a $1$ brick wall, or a $1$ brick wall and a $1$ brick wall. The nim-values of these states are $2$, $3$, $3$, and $0$, respectively, and hence the nim-value of this game state is $1$.

 [Why is the nim-value of it $1$? - awesomediabrine

 \href{https://en.wikipedia.org/wiki/Mex\_(mathematics)}{https://en.wikipedia.org/wiki/Mex\_(mathematics)}]

 The possible game states after the $5$ brick wall are the following: a $3$ brick wall, a $4$ brick wall, a $3$ brick wall and a $1$ brick wall, a two $2$ brick walls, and a $2$ brick wall plus a $1$ brick wall. The nim-values of these are $3$, $1$, $2$, $0$, and $3$, respectively, meaning the nim-value of a $5$ brick wall is $4$. 

 Finally, we find the nim-value of a $6$ brick wall. The possible states are a $5$ brick wall, a $4$ brick wall and a $1$ brick wall, a $3$ brick wall and a $2$ brick wall, a $4$ brick wall, a $3$ brick wall and a $1$ brick wall, and finally two $2$ brick walls. The nim-values of these game states are $4$, $0$, $1$, $1$, $2$, and $0$, respectively. This means the $6$ brick wall has a nim-value of $3$. 

 The problem is asking which of the answer choices is losing, or has a nim-value of $0$. We see that option $\textbf{(A)}$ has a nim-value of $3\oplus1\oplus1=3$, option $\textbf{(B)}$ has a nim-value of $3\oplus2\oplus1=0$, option $\textbf{(C)}$ has a nim-value of $3\oplus2\oplus2=3$, option $\textbf{(D)}$ has a nim-value of $3\oplus3\oplus1=1$, and option $\textbf{(E)}$ has a nim-value of $3\oplus3\oplus2=2$, so the answer is $\boxed{\textbf{(B) }(6, 2, 1)}$.

 This method can also be extended to solve the note after the first solution. The nim-values of the $7$ brick wall and the $8$ brick wall are $2$ and $1$, using the same method as above. The nim-value of $(8, 7, 5, 3, 2)$ is therefore $1\oplus2\oplus4\oplus3\oplus2 = 6$, which is winning.

 



\textbf{Solution 4}

Consider the following, much simpler game.

 Arjun and Beth can each either take 1 or 2 bricks from the right-hand-side of a continuous row of initially $n$ bricks.It is easy to see that for $n$ a multiple of 3, Beth can "mirror" whatever Arjun plays: if he takes 1, she takes 2, and if he takes 2, she takes 1. With this strategy, Beth always takes the last brick. If $n$ is not a multiple of 3, then Arjun takes whichever amount puts Beth in the losing position.

 The total number of bricks in the initial states given by the answer choices is $8, 9, 10, 10, 11$. Thus, answer choice $\textbf{(B)}$ appears promising as a winning position for Beth. The difference between this game and the simplified game is that in certain positions, namely those consisting of fragments of size only 1, taking 2 bricks is not allowed. We can assume that for the starting position $(6, 2, 1)$, Beth always has a move to ensure that she can continue to mirror Arjun throughout the game. (This could be proven rigorously with lots of casework. In particular, she must avoid providing Arjun with a position of 3 continuous bricks, because then he could take the middle block and force a win.) The assumption seems reasonable, so the answer is $\boxed{\textbf{(B)}}$.

 



\textbf{Solution 5}

We can start by guessing and checking our solutions. Let's start with $A$, $(6,1,1)$. Arjun goes first, and he has a winning strategy. His strategy is to cut the 6 block in half, so whatever Beth does, Arjun will copy. If we see this in practice, Arjun has made the block in to $(2,2,1,1)$. If Beth takes 1, he will take one. If Beth takes 2, he will take 2 as well.

 Then, we can guess solution $B$, $(6,2,1)$. If Arjun starts by taking 1 of the blocks, he will either have $(6,1,1)$, $(6,2)$, $(x,y,2,1)$ (Where $x+y = 5$). If he has $(6,1,1)$, Beth can take 2 out of the middle of $(6,1,1)$ to get $(2,2,1,1)$ which we know Arjun will lose because we know by symmetry from answer choice $A$. If he has $(6,2)$, Beth will take out one from the edge, and we get $(5,2)$. Beth can now copy whatever Arjun does, so she will win. 

 Next, if he takes one out to get $(x,y,2,1)$, we can either have $(4,1,2,1)$ or $(3,2,2,1)$. If it is the first case, Beth can take the edge 2 out of the 4, which gives us $(2,1,2,1)$ and whatever Arjun does, Beth can do, so she has established a win. If it is the 2nd case, Beth can take out 2 from the 3 to get $(1,2,2,1)$, in which, Beth can copy whatever Arjun does, so Beth will win. So in all 4 cases, Beth wins. So our answer is  $\boxed{\textbf{(B)}}$




\textbf{Solution 6 (Mirroring)}

In many game problems, the winning strategy is to simply copy all of your opponent's moves. Since you will always have a valid move if you mirror a valid move, you must move last and win.

 Then, if Arjun can make the game "mirrorable" in his first move, then he can just copy all of Beth's moves after that.

 $(6,1,1)$ can be turned into $(2,2,1,1)$ by Arjun. Any move Beth makes, Arjun can just copy on the other block with same value. For example, if Beth took one block off the first $2$, Arjun takes one block off the second $2$.

 $(6,2,2)$ turns into $(2,2,2,2)$.

 $(6,3,1)$ turns into $(3,1,3,1)$.

 $(6,3,2)$ turns into $(3,2,3,2)$.

 There are definitely other ways to split, but Arjun must win these games. This leaves answer $\boxed{B(6,2,1)}$.

 



\textbf{Solution 7 (Side Games)}

We will call the $6$-wall the main game and the remaining walls the side game.It is clear that if there is no side game, Arjun can win by taking the middle two walls out. However, if Beth is ever able to force Arjun to take two turns in a row on the main game, their roles will swap and she will win. To do so, she must force Arjun to take two turns in a row, which means she needs to take the last wall in the side game. In other words, we can reduce the problem into Beth winning a game with starting configuration ignoring the $6$-wall.

 From here, it's clear that Beth loses $(1,1)$, $(2,2)$, and $(3,2)$. However, Beth does win both $(3,1)$ and $(2,1)$. However, with the $(3,1)$ case, Arjun can instead start the game by splitting the $(3,1)$ into $(1, 1, 1)$, which would swap their roles. Since Arjun now wins the side game, Beth will lose the $(3,1)$ case, making the answer $\boxed{B(6, 2, 1)}$



	\item (\textit{AMC 12 12B 2021}) Three balls are randomly and independently tossed into bins numbered with the positive integers so that for each ball, the probability that it is tossed into bin $i$ is $2^{-i}$ for $i=1,2,3,....$ More than one ball is allowed in each bin. The probability that the balls end up evenly spaced in distinct bins is $\frac pq,$ where $p$ and $q$ are relatively prime positive integers. (For example, the balls are evenly spaced if they are tossed into bins $3,17,$ and $10.$) What is $p+q?$
$\textbf{(A) }55 \qquad \textbf{(B) }56 \qquad \textbf{(C) }57\qquad \textbf{(D) }58 \qquad \textbf{(E) }59$

 \textbf{Answer: }A



\textbf{Solution 1}

"Evenly spaced" just means the bins form an arithmetic sequence.

 Suppose the middle bin in the sequence is $x$. There are $x-1$ different possibilities for the first bin, and these two bins uniquely determine the final bin. Now, the probability that these $3$ bins are chosen is $6\cdot 2^{-3x} = 6\cdot \frac{1}{8^x}$, so the probability $x$ is the middle bin is $6\cdot\frac{x-1}{8^x}$. Then, we want the sum
 \begin{align*} 6\sum_{x=2}^{\infty}\frac{x-1}{8^x} &= \frac{6}{8}\left[\frac{1}{8} + \frac{2}{8^2} + \frac{3}{8^3}\cdots\right]\\ &= \frac34\left[\left(\frac{1}{8} + \frac{1}{8^2} + \frac{1}{8^3}+\cdots \right) + \left(\frac{1}{8^2} + \frac{1}{8^3} + \frac{1}{8^4} + \cdots \right) + \cdots\right]\\ &= \frac34\left[\frac17\cdot \left(1 + \frac{1}{8} + \frac{1}{8^2} + \frac{1}{8^3} + \cdots \right)\right]\\ &= \frac34\cdot \frac{8}{49}\\ &= \frac{6}{49} \end{align*}
The answer is $6+49=\boxed{\textbf{(A) }55}.$




\textbf{Solution 2}

As in solution 1, note that "evenly spaced" means the bins are in arithmetic sequence. We let the first bin be $a$ and the common difference be $d$. Further note that each $(a, d)$ pair uniquely determines a set of $3$ bins. 

 We have $a\geq1$ because the leftmost bin in the sequence can be any bin, and $d\geq1$, because the bins must be distinct.

 This gives us the following sum for the probability:
 \begin{align*} 6 \sum_{a=1}^{\infty} \sum_{d=1}^{\infty} 2^{-3a-3d} &= 6 \sum_{a=1}^{\infty} \sum_{d=1}^{\infty} 2^{-3a} \cdot 2^{-3d} \\ &= 6 \left( \sum_{a=1}^{\infty} 2^{-3a} \right) \left( \sum_{d=1}^{\infty} 2^{-3d} \right) \\ &= 6 \left( \sum_{a=1}^{\infty} 8^{-a} \right) \left( \sum_{d=1}^{\infty} 8^{-d} \right) \\ &= 6 \left( \frac{1}{7} \right) \left( \frac{1}{7} \right) \\ &= \frac{6}{49} .\end{align*}
Therefore the answer is $6 + 49 = \boxed{\textbf{(A) }55}$.

 -Darren Yao

 



\textbf{Solution 3}

This is a slightly messier variant of solution 2. If the first ball is in bin $i$ and the second ball is in bin $j>i$, then the third ball is in bin $2j-i$. Thus the probability is
 \begin{align*} 6\sum_{i=1}^{\infty}\sum_{j=i+1}^\infty2^{-i}2^{-j}2^{-2j+i}&=6\sum_{i=1}^{\infty}\sum_{j=i+1}^\infty2^{-3j}\\ &=6\sum_{i=1}^{\infty}\left(\frac{2^{-3(i+1)}}{1-\frac{1}{8}}\right)\\ &=6\sum_{i=1}^\infty\frac{8}{7}\cdot2^{-3}\cdot2^{-3i}\\ &=\frac{6}{7}\sum_{i=1}^\infty2^{-3i}\\ &=\frac{6}{7}\cdot\frac{2^{-3}}{1-\frac18}\\ &=\frac{6}{49}. \end{align*}
Therefore the answer is $6 + 49 = \boxed{\textbf{(A) }55}$.

 



\textbf{Solution 4 (Table)}

Based on the value of $n,$ we construct the following table:
 \[\begin{array}{c|c|c|c}  & & & \\ [-1.5ex] \textbf{Exactly }\boldsymbol{n}\textbf{ Spaces Apart} & \textbf{Bin \#s} & \textbf{Expression} & \textbf{Prob. of One Such Perm.} \\ [1ex]  \hline\hline   & & & \\ [-1.5ex]  n=1 & 1,2,3 & 2^{-1}\cdot2^{-2}\cdot2^{-3} & 2^{-6} \\ [1ex]   & 2,3,4 & 2^{-2}\cdot2^{-3}\cdot2^{-4} & 2^{-9} \\ [1ex]   & 3,4,5 & 2^{-3}\cdot2^{-4}\cdot2^{-5} & 2^{-12} \\ [1ex]   & 4,5,6 & 2^{-4}\cdot2^{-5}\cdot2^{-6} & 2^{-15} \\ [1ex]   & \cdots & \cdots & \cdots \\ [1ex]  \hline  & & & \\ [-1.5ex]  n=2 & 1,3,5 & 2^{-1}\cdot2^{-3}\cdot2^{-5} & 2^{-9} \\ [1ex]   & 2,4,6 & 2^{-2}\cdot2^{-4}\cdot2^{-6} & 2^{-12} \\ [1ex]   & 3,5,7 & 2^{-3}\cdot2^{-5}\cdot2^{-7} & 2^{-15} \\ [1ex]   & 4,6,8 & 2^{-4}\cdot2^{-6}\cdot2^{-8} & 2^{-18} \\ [1ex]   & \cdots & \cdots & \cdots \\ [1ex]  \hline  & & & \\ [-1.5ex]  n=3 & 1,4,7 & 2^{-1}\cdot2^{-4}\cdot2^{-7} & 2^{-12} \\ [1ex]   & 2,5,8 & 2^{-2}\cdot2^{-5}\cdot2^{-8} & 2^{-15} \\ [1ex]   & 3,6,9 & 2^{-3}\cdot2^{-6}\cdot2^{-9} & 2^{-18} \\ [1ex]  & 4,7,10 & 2^{-4}\cdot2^{-7}\cdot2^{-10} & 2^{-21} \\ [1ex]  & \cdots & \cdots & \cdots \\ [1ex]  \hline   & & & \\ [-1.5ex]    \cdots & \cdots & \cdots & \cdots \\ [1ex] \end{array}\]Since three balls have $3!=6$ permutations, the requested probability is
 \begin{align*} 6\left(\sum_{k=0}^{\infty}2^{-6-3k}+\sum_{k=0}^{\infty}2^{-9-3k}+\sum_{k=0}^{\infty}2^{-12-3k}+\cdots\right)&=6\left(2^{-6}\sum_{k=0}^{\infty}2^{-3k}+2^{-9}\sum_{k=0}^{\infty}2^{-3k}+2^{-12}\sum_{k=0}^{\infty}2^{-3k}+\cdots\right) \\ &=\left(6\sum_{k=0}^{\infty}2^{-3k}\right)\cdot\left(2^{-6}+2^{-9}+2^{-12}+\cdots\right) \\ &=\left(6\sum_{k=0}^{\infty}2^{-3k}\right)\cdot\left(\sum_{k=0}^{\infty}2^{-6-3k}\right) \\ &=\left(6\sum_{k=0}^{\infty}2^{-3k}\right)\cdot\left(2^{-6}\sum_{k=0}^{\infty}2^{-3k}\right) \\ &=\frac{6}{1-2^{-3}}\cdot\frac{2^{-6}}{1-2^{-3}} \\ &=\frac{6}{49} \end{align*}
by infinite geometric series, from which the answer is $6+49=\boxed{\textbf{(A) }55}.$



	\item (\textit{AMC 12 12B 2021}) Let $ABCD$ be a parallelogram with area $15$. Points $P$ and $Q$ are the projections of $A$ and $C,$ respectively, onto the line $BD;$ and points $R$ and $S$ are the projections of $B$ and $D,$ respectively, onto the line $AC.$ See the figure, which also shows the relative locations of these points.


 
                    \begin{figure}[H]
                        \centering
                        \includegraphics[width=0.5\textwidth]{Problem/AMC_12_12B_2021_24_Problem_Figure_1.png}
                    \end{figure}
                    


 Suppose $PQ=6$ and $RS=8,$ and let $d$ denote the length of $\overline{BD},$ the longer diagonal of $ABCD.$ Then $d^2$ can be written in the form $m+n\sqrt p,$ where $m,n,$ and $p$ are positive integers and $p$ is not divisible by the square of any prime. What is $m+n+p?$
$\textbf{(A) }81 \qquad \textbf{(B) }89 \qquad \textbf{(C) }97\qquad \textbf{(D) }105 \qquad \textbf{(E) }113$

 \textbf{Answer: }A



\textbf{Solution 1}

Let $X$ denote the intersection point of the diagonals $AC$ and $BD$. Remark that by symmetry $X$ is the midpoint of both $\overline{PQ}$ and $\overline{RS}$, so $XP = XQ = 3$ and $XR = XS = 4$. Now note that since $\angle APB = \angle ARB = 90^\circ$, quadrilateral $ARPB$ is cyclic, and so
 \[XR\cdot XA = XP\cdot XB,\]which implies $\frac{XA}{XB} = \frac{XP}{XR} = \frac34$.

 Thus let $x> 0$ be such that $XA = 3x$ and $XB = 4x$. Then Pythagorean Theorem on $\triangle APX$ yields $AP = \sqrt{AX^2 - XP^2} = 3\sqrt{x^2-1}$, and so
 \[[ABCD] = 2[ABD] = AP\cdot BD = 3\sqrt{x^2-1}\cdot 8x = 24x\sqrt{x^2-1}=15\]Solving this for $x^2$ yields $x^2 = \frac12 + \frac{\sqrt{41}}8$, and so
 \[(8x)^2 = 64x^2 = 64\left(\frac12 + \frac{\sqrt{41}}8\right) = 32 + 8\sqrt{41}.\]The requested answer is $32 + 8 + 41 = \boxed{\textbf{(A)} ~81}$.

 



\textbf{Solution 2 (Trig)}

Let $X$ denote the intersection point of the diagonals $AC$ and $BD,$ and let $\theta = \angle{COB}$. Then, by the given conditions, $XR = 4,$ $XQ = 3,$ $[XCB] = \frac{15}{4}$. So,
 \[XC = \frac{3}{\cos \theta}\]
 \[XB \cos \theta = 4\]
 \[\frac{1}{2} XB XC \sin \theta = \frac{15}{4}\]Combining the above 3 equations, we get
 \[\frac{\sin \theta }{\cos^2 \theta} = \frac{5}{8}.\]Since we want to find $d^2 = 4XB^2 =  \frac{64}{\cos^2 \theta},$ we let $x = \frac{1}{\cos^2 \theta}.$ Then 
 \[\frac{\sin^2 \theta }{\cos^4 \theta} = \frac{1-\cos ^2 \theta}{\cos^4 \theta} = x^2 - x = \frac{25}{64}.\]Solving this, we get $x = \frac{4 + \sqrt{41}}{8},$ so $d^2 = 64x = 32 + 8\sqrt{41} \rightarrow 32+8+41=\boxed{\textbf{(A)} ~81}$




\textbf{Solution 3 (Similar Triangles and Algebra)}

Let $X$ be the intersection of diagonals $AC$ and $BD$. By symmetry $[\triangle XCB] = \frac{15}{4}$, $XQ = 3$ and $XR = 4$, so now we have reduced all of the conditions one quadrant. Let $CQ = x$. $XC = \sqrt{x^2+9}$, $RB = \frac{4x}{3}$ by similar triangles and using the area condition we get $\frac{4}{3} \cdot x \cdot \sqrt{x^2+9} = \frac{15}{2}$. Note that it suffices to find $XB = \frac{4}{3}\sqrt{x^2+9}$ because we can double and square it to get $d^2$. Solving for $a = x^2$ in the above equation, and then using $d^2 = \frac{64}{9}(x^2+9) = 8\sqrt{41} + 32 \Rightarrow 8+41+32=\boxed{\textbf{(A)} ~81}$




\textbf{Solution 4 (Similar Triangles)}

Again, Let $X$ be the intersection of diagonals $AC$ and $BD$. Note that triangles $\triangle QXC$ and $\triangle BXR$ are similar because they are right triangles and share $\angle CXQ$. First, call the length of $XB = \frac{d}{2}$. By the definition of an area of a parallelogram, $CQ \cdot 2XB = 15$, so $CQ = \frac{15}{d}$. Using similar triangles on $\triangle QXC$ and $\triangle BXR$, $\frac{CQ}{XQ} = \frac{BR}{XR}$. Therefore, finding $BR$, $BR = \frac{XR}{XQ} \cdot CQ = \frac{4}{3} \cdot \frac{15}{d} = \frac{20}{d}$. Now, applying the Pythagorean theorem once, we find $(\frac{20}{d}) ^2$ + $(4)^2$ = $(\frac{d}{2}) ^2$. Solving this equation for $d^2$, we find $d^2=\frac{64+\sqrt{4096+6400}}{2}=32+8\sqrt{41} \rightarrow 32+8+41= \boxed{\textbf{(A)} ~81}$




\textbf{Solution 5}

Let $BQ = PD = x.$ We know that the area of the parallelogram is $15,$ so it follows that $[\triangle{BCD}] = [\triangle{BAD}] = \frac{15}{2}$ and the height of each triangle, which are also the lengths of $QC$ and $AP,$ is $\frac{15}{2(x+3)}.$ Suppose that $E = RS \cap BD.$ Because $\angle{BRE} = \angle{CQE}$ and $\angle{BER} = \angle{CQD},$ we have $\triangle{BRE} \sim \triangle{CQE}.$ The length of $CE,$ by the Pythagorean Theorem is $\sqrt{3^2+(\frac{15}{2(x+3)})^2}$ and the length of $BR,$ by the Pythagorean Theorem on $\triangle{BRE},$ is $\sqrt{(x+3)^2 - 4}.$ Note that
 \[\sin{\angle QEC} = \frac{CQ}{CE} = \frac{BR}{BE}\]Substituting in our values,
 \[\frac{\frac{15}{2(x+3)}}{\sqrt{9+(\frac{15}{2(x+3)})^2}} = \frac{\sqrt{(x+3)^2 - 4^2}}{x+3}\]To rid unnecessary computation, we let $(x+3)^2 = a.$ The equation simplifies, after cross multiplying, to
 \[\sqrt{9+\frac{15^2}{4a}} \sqrt{a-16 } = \frac{15}{2}\]
 \[36a^2 - 576a - 15^2\cdot 16 = 0\]
 \[a^2-16a-100 =0\]By the quadratic formula, $a \in \{\frac{16 - \sqrt{656}}{2}, \frac{16 + \sqrt{656}}{2}\},$ so we discard the negative solution. The value of $BD^2$ is 
 \[BD^2 = (2x+6)^2 = 4(x+3)^2 = 4a = 4 \cdot \frac{16 + \sqrt{656}}{2} = 32+8\sqrt{41}\] and the desired answer is $32+8+41 = \boxed{\textbf{(A)} ~81}$ ~skyscraper

 



\textbf{Solution 6 (Similar Triangles & Pythagorean Theorem)}

Let the intersection of $RS$ and $BD$ be $X$. 

 $\because$ $\angle APX = \angle DSX$ and $\angle AXP = \angle DXS$, $\triangle APX \sim \triangle DSX$ by $AA$
$\therefore$ $\frac{PA}{DS} = \frac68 = \frac34$, $DS = \frac43 \cdot PA$

 By the Pythagorean theorem and the property of projection, $BD^2 = (DS+BR)^2 + RS^2 = 4DS^2 + 64 = 4(\frac43 \cdot PA)^2 + 64 = \frac{64}{9} \cdot PA^2 + 64$, $\frac{64}{9} \cdot PA^2 = BD^2 - 64$
$\because [ABCD] = PA \cdot BD = 15$, $\therefore PA = \frac{15}{BD}$

 \[\frac{64}{9} (\frac{15}{BD})^2 = BD^2 - 64\]

 \[\frac{1600}{BD^2} = BD^2 - 64\]

 \[BD^4 - 64 BD^2 - 1600 = 0\]

 \[BD^2 = \frac{64 + \sqrt{64^2 - 4 (-1600)}}{2} = 32 + 8 \sqrt{41}\]

 Therefore, the answer is $32 + 8 + 41 = \boxed{\textbf{(A) } 81}$.

 



\textbf{Solution 7}

Let $O$ be the intersection of the diagonals in quadrilateral $ABCD$ and let the origin be at $O$.

 Then, let $BD$ be on the x-axis, making $\vec{Q}=\begin{pmatrix}3 \\ 0\end{pmatrix}$, $\vec{C}=\begin{pmatrix}3 \\ y\end{pmatrix}$, $\vec{B}=\begin{pmatrix}x \\ 0\end{pmatrix}$, and $\vec{D}=\begin{pmatrix}-x \\ 0\end{pmatrix}$.

 The area of $\triangle BDC$ is half the area of $ABCD$, so,
 \[\frac{1}{2}(\overline{BD})(\overline{QC})=\frac{1}{2}(ABCD)\]
 \[\frac{1}{2}(2x)(y)=\frac{1}{2}(15)\]
 \[xy=\frac{15}{2}\]
 \[y= \frac{15}{2x}\].

 Also, $\vec{R}$ is the projection of $\vec{B}$ onto $\vec{AC}$, which is the same as projecting it onto $\vec{C}$. We get:
 \[\vec{R}=proj_{\vec{C}}{(\vec{B})}=\frac{\vec{B} \cdot \vec{C}}{\vec{C} \cdot \vec{C}}(\vec{C})=\frac{\begin{pmatrix}x \\ 0\end{pmatrix} \cdot \begin{pmatrix}3 \\ y\end{pmatrix}}{\begin{pmatrix}3 \\ y\end{pmatrix} \cdot \begin{pmatrix}3 \\ y\end{pmatrix}}\begin{pmatrix}3 \\ y\end{pmatrix}=\frac{3x}{y^2+9}\begin{pmatrix}3 \\ y\end{pmatrix}\]

 We are given that $\overline{RS} = 8$, so $||\vec{RS}|| = 8$. Because diagonals bisect each other in a parallelogram, $||\vec{R}|| = \frac{8}{2} = 4$. Substituting $||\vec{R}||$ with the previous equation gives:
 \[||\frac{3x}{y^2+9}\begin{pmatrix}3 \\ y\end{pmatrix}|| = 4\]
 \[\frac{3x}{y^2+9}\sqrt{y^2+9} = 4\]
 \[\frac{3x}{\sqrt{y^2+9}} = 4\]Squaring both sides gives:
 \[\frac{9x^2}{y^2+9}=16\]
 \[9x^2=16y^2+144\]Substituting $y=\frac{15}{2x}$:
 \[9x^2=16\left(\frac{15}{2x}\right)^2+144\]
 \[9x^2=16\left(\frac{225}{4x^2}\right)+144\]
 \[9x^4=144x^2+900\]
 \[x^4-16x^2-100=0\]Solving the quadratic for $x^2$ gives:
 \[x^2 = \frac{16 \pm \sqrt{16^2 - 4(1)(-100)}}{2} = 8 \pm 2\sqrt{41}\]But $x^2$ is nonnegative, so $x^2$ must be $8 + 2\sqrt{41}$. We are looking for $\overline{BD}^2$, which is $4x^2$. $4(8 + 2\sqrt{41}) = 32 + 8\sqrt{41}$, so $a + b + c = 32 + 8 + 41 = \boxed{\textbf{(A) } 81}$ ~askeww (My first solution ever!)

 



\textbf{Solution 8}

First and foremost, the condition given in the problem could be represented. Moreover, let $QB=x$ and $CQ=h$ for convenience.
                    \begin{figure}[H]
                        \centering
                        \includegraphics[width=0.5\textwidth]{Solution/AMC_12_12B_2021_24_Solution_Figure_1.png}
                    \end{figure}
                    Using the property of a parallelogram, it could be inferred that $\frac{h(3+x)}{2}=\frac{15}{4}$. Moreover, Pythagorean Theorem manifests that $\left(\frac{4}{3}h\right)^2+4^2=(3+x)^2$.
 \begin{align*}     3+x&=\frac{15}{2h} \\     \left(\frac{4}{3}h\right)^2+4^2&=\left(\frac{15}{2h}\right)^2 \\     \frac{16h^2}{9}+16&=\frac{225}{4h^2} \\     64h^4+576h^2-2025&=0 \\     h^2&\Rightarrow\frac{9\sqrt{41}-36}{8} \end{align*}

 \begin{align*}     \left(\frac{4}{3}h\right)^2+4^2&=(3+x)^2 \\     \frac{16}{9}\cdot\frac{9\sqrt{41}-36}{8}+4^2&=(3+x)^2 \\     2\sqrt{41}-8+16&=(3+x)^2 \\     (3+x)^2&=8+2\sqrt{41} \\     d^2=4(3+x)^2&=32+8\sqrt{41} \\ \end{align*}
Thus, $32+8+41=\boxed{\textbf{(A) }81}$.

 



\textbf{Solution 9(Similar to other solutions)}

We first note that $DS$ and $BR$ are perpendicular to $AC$. Then, let $DS = BR = h$. Furthermore, let $AS = x$ and the intersections of diagonals $AC$ and $BD$ be $T$. We note that triangle $APT$ is similar to triangle $STD$. We will also show that the areas of triangles $ADT, BCT, CDT, ABT$ are the same namely $\frac{15}{4}$. Let $\angle DTA = \theta$. Then, $\angle DTC = \pi - \angle DTA = \pi - \theta$. Let $DP = QB = y$. Obviously, $PT = QT = \frac{6}{2} = 3$. Also, $AT = CT = 4 - x$. Then the area of triangle $ADT$ is just $\frac{1}{2} \cdot (y + 3)(4 - x) \cdot \sin(\theta)$. It is easy to see that $CBT$ is congruent to $ADT$ and $ABT$ is congruent to $CDT$ so it suffices to show triangle $CDT$ has the same area as triangle $ADT$. In fact, the area equals $\frac{1}{2} \cdot (y + 3)(4 - x) \cdot \sin(\pi - \theta) = \frac{1}{2} \cdot (y + 3)(4 - x) \cdot \sin(\theta)$ which indeed is equal to the area of triangle $ADT$. Hence, each of these triangles breaks the area of the entire parallelogram into $4$ pieces so each area is $\frac{15}{4}$. Look at triangle $ADT$. The base is $y + 3$. Let the height equal $m$. Hence:

 $\frac{1}{2} \cdot m \cdot (y + 3) = \frac{15}{4} \implies m = \frac{15}{2(y + 3)}$

 Remember that we established the fact that triangle $APT$ is similar to triangle $STD$ and so:

 $\frac{4}{3} = \frac{\frac{15}{2(y + 3)}}{h} \implies h(y + 3) = 10$

 Also by the Pythagorean Theorem on triangle $STD$
$h^{2} + 16 = (y + 3)^{2} \implies h^{2} + 16 = \frac{100}{h^{2}} \implies h^{4} + 16h^{2} - 100 = 0 \implies h^{2} = \frac{-16 \pm \sqrt{16 \cdot 16 - 4(-100)}}{2} = \frac{-16 \pm \sqrt{16 \cdot 16 + 16 \cdot 25}}{2} = \frac{-16 \pm 4\sqrt{41}}{2} = -8 \pm 2\sqrt{41}$.

 Since $h^{2}$ is positive, $h^{2} = -8 + 2\sqrt{41}$. We need to find $BD^{2} = (y + 3 + 3 + y)^{2} = 4(y + 3)^{2}$ so remember $(y + 3)^{2} = \frac{100}{h^{2}} = \frac{100}{-8 + 2\sqrt{41}} = 8 + 2\sqrt{41}$. But $BD^{2}$ is $4$ times this so the answer is $4(8 + 2\sqrt{41}) = 32 + 8\sqrt{41}$. Hence $32 + 8 + 41 = 81 \implies \boxed{\textbf{(A)}}$



	\item (\textit{AMC 12 12B 2021}) Let $S$ be the set of lattice points in the coordinate plane, both of whose coordinates are integers between $1$ and $30,$ inclusive. Exactly $300$ points in $S$ lie on or below a line with equation $y=mx.$ The possible values of $m$ lie in an interval of length $\frac ab,$ where $a$ and $b$ are relatively prime positive integers. What is $a+b?$
$\textbf{(A) }31 \qquad \textbf{(B) }47 \qquad \textbf{(C) }62\qquad \textbf{(D) }72 \qquad \textbf{(E) }85$

 \textbf{Answer: }E



\textbf{Solution 1}

First, we find a numerical representation for the number of lattice points in $S$ that are under the line $y=mx.$ For any value of $x,$ the highest lattice point under $y=mx$ is $\lfloor mx \rfloor.$ Because every lattice point from $(x, 1)$ to $(x, \lfloor mx \rfloor)$ is under the line, the total number of lattice points under the line is $\sum_{x=1}^{30}(\lfloor mx \rfloor).$

 
Now, we proceed by finding lower and upper bounds for $m.$ To find the lower bound, we start with an approximation. If $300$ lattice points are below the line, then around $\frac{1}{3}$ of the area formed by $S$ is under the line. By using the formula for a triangle's area, we find that when $x=30, y \approx 20.$ Solving for $m$ assuming that $(30, 20)$ is a point on the line, we get $m = \frac{2}{3}.$ Plugging in $m$ to $\sum_{x=1}^{30}(\lfloor mx \rfloor),$ we get:

 
 \[\sum_{x=1}^{30}(\lfloor \frac{2}{3}x \rfloor) = 0 + 1 + 2 + 2 + 3 + \cdots + 18 + 18 + 19 + 20\]

 We have a repeat every $3$ values (every time $y=\frac{2}{3}x$ goes through a lattice point). Thus, we can use arithmetic sequences to calculate the value above:

 
 \[\sum_{x=1}^{30}(\lfloor \frac{2}{3}x \rfloor) = 0 + 1 + 2 + 2 + 3 + \cdots + 18 + 18 + 19 + 20\]
 \[=\frac{20(21)}{2} + 2 + 4 + 6 + \cdots + 18\]
 \[=210 + \frac{20}{2}\cdot 9\]
 \[=300\]

 This means that $\frac{2}{3}$ is a possible value of $m.$ Furthermore, it is the lower bound for $m.$ This is because $y=\frac{2}{3}x$ goes through many points (such as $(21, 14)$). If $m$ was lower, $y=mx$ would no longer go through some of these points, and there would be less than $300$ lattice points under it.

 Now, we find an upper bound for $m.$ Imagine increasing $m$ slowly and rotating the line $y=mx,$ starting from the lower bound of $m=\frac{2}{3}.$The upper bound for $m$ occurs when $y=mx$ intersects a lattice point again (look at this problem to get a better idea of what's happening: \href{https://artofproblemsolving.com/wiki/index.php/2011\_AMC\_10B\_Problems/Problem\_24}{https://artofproblemsolving.com/wiki/index.php/2011\_AMC\_10B\_Problems/Problem\_24}).

 In other words, we are looking for the first $m > \frac{2}{3}$ that is expressible as a ratio of positive integers $\frac{p}{q}$ with $q \le 30.$ For each $q=1,\dots,30$, the smallest multiple of $\frac{1}{q}$ which exceeds $\frac{2}{3}$ is $1, \frac{2}{2}, \frac{3}{3}, \frac{3}{4}, \frac{4}{5}, \cdots , \frac{19}{27}, \frac{19}{28}, \frac{20}{29}, \frac{21}{30}$ respectively, and the smallest of these is $\frac{19}{28}.$

 Alternatively, see the method of finding upper bounds in solution 2.

 The lower bound is $\frac{2}{3}$ and the upper bound is $\frac{19}{28}.$ Their difference is $\frac{1}{84},$ so the answer is $1 + 84 = \boxed{85}.$

 An alternative would be using Farey fractions and the mediant theorem to find the upper bound. $\frac{2}{3}$ and $\frac{7}{10}$ gives $\frac{9}{13},$ and so on using Farey addition. 

 The above alternative can be made quicker by using a different Farey sequence rule. The Farey sequence rule states that if $\frac{a}{b}$ and $\frac{c}{d}$ are adjacent in a Farey sequence, we have bc - ad = 1 (see \href{https://en.wikipedia.org/wiki/Farey\_sequence\#Farey\_neighbours}{https://en.wikipedia.org/wiki/Farey\_sequence\#Farey\_neighbours}). We use the rule with a = 2, b = 3, so we have 3c - 2d = 1. Starting with d = 30 (since d $\leq$ 30), we see that c would be noninteger and the same can be shown for d = 29. When d = 28, c = 19 so the range is [$\frac{2}{3}$, $\frac{19}{28}$] => $\frac{1}{84}$ => $\boxed{85}$

 
An alternative approach with the same methodology can be done using Pick's Theorem. Wikipedia page: \href{https://en.wikipedia.org/wiki/Pick\%27s\_theorem}{https://en.wikipedia.org/wiki/Pick\%27s\_theorem}. It's a formula to find the number of lattice points strictly inside a polygon. Approximation of the lower bound is still necessary.

 



\textbf{Solution 2}

Based on area ratios between a square of side length $30$ and a triangle with base $30$, we estimate that the slope $m$ of the line we want is approximately $\frac{2}{3}$. Following this estimate, we see if there are approximately $30 \cdot 30 - 300 = 600$ lattice points above the line $y=\frac{2}{3}x$.

 Counting the number of lattice points with $x=1$ above the line, the number of lattice points with $x=2$ above the line, and so on, we find that the total number of lattice points above the line is $30+29+28+28+27+26+26 \ldots+ 10$, with the even integers repeating every third term. We see that the average of the $30$ terms is $20$, which means that exactly $20 \cdot 30 = 600$ lattice points above the line as desired. This gives a lower bound because any decrease in the slope of the line would cause points that were already on the line to shift to being above it.

 To find the upper bound, notice that each lattice point less than $1$ unit above the line is either $\frac{1}{3}$ or $\frac{2}{3}$ above. Since the slope through a point is the y-coordinate divided by the x-coordinate, a shift in the slope will increase the y-value of the higher x-coordinates more than the y-value of the lower x-coordinates. So, we turn our attention to $x=28, 29, 30$ for which the line $y=\frac{2}{3}x$ intersects at $y= \frac{56}{3}, \frac{58}{3}, 20$. The point $(30,20)$ is already counted, and we can clearly see that if we slowly increase the slope of the line, we will first hit the point $(28,19)$ since $(28, \frac{56}{3})$ is the closest to it. The equation of the line which goes through both the origin and $(28,19)$ is $y=\frac{19}{28}x$. This gives an upper bound of $m=\frac{19}{28}$. 

 Taking the upper bound of $m$ and subtracting the lower bound yields $\frac{19}{28}-\frac{2}{3}=\frac{1}{84}$. The answer is therefore $1+84=$ $\boxed{\textbf{(E)} ~85}$.

 \textbf{Diagram} 
                    \begin{figure}[H]
                        \centering
                        \includegraphics[width=0.5\textwidth]{Solution/AMC_12_12B_2021_25_Solution_Figure_1.png}
                    \end{figure}
                    




\textbf{Solution 3}

As the procedure shown in the Solution 1, the lower bound of $m$ is $\frac{2}{3}.$Here I give a more logic way to show how to find the upper bound of $m.$Denote $N=\sum_{x=1}^{30}(\lfloor mx \rfloor)$ as the number of lattice points in $S$.

 $N = \lfloor m \rfloor+\lfloor 2m \rfloor+\lfloor 3m \rfloor+\cdots+\lfloor 30m \rfloor = 300 .$

 Let $m = \frac{2}{3}+k$. for $\forall x_{i}\le 30, x\in N^{*},  \lfloor mx_{i} \rfloor = \lfloor \frac{2}{3}x+xk \rfloor.$

 Our target is finding the minimum value of $k$ which can increase one unit of $\lfloor mx_{i} \rfloor .$

 Notice that:

 When $x_{i} = 3n, \lfloor mx_{i} \rfloor = \lfloor 2n+(3n)k \rfloor$ We don't have to discuss this case.

 When $x_{i} = 3n+1, \lfloor mx_{i} \rfloor = \lfloor 2n+\frac{2}{3}+(3n+1)k \rfloor, k_{min1}=\frac{1}{3(3n+1)}.$

 When $x_{i} = 3n+2, \lfloor mx_{i} \rfloor = \lfloor 2n+1+\frac{1}{3}+(3n+2)k \rfloor, k_{min2}=\frac{2}{3(3n+2)}.$ 

 Here $n\in N^{*}, n \le 9.$

 Denote $k_{min}=min\left \{k_{min1},k_{min2}  \right \}.$

 Obviously $k_{min1}$ and $k_{min2}$ are decreasing. Let's considering the situation when $n=9.$
$k_{min}=min\left\{\frac{1}{84},\frac{2}{87}\right\}=\frac{1}{84}.$

 This means that the answer is just $\frac{1}{84}$, so $a+b=85$.  Choose $\boxed{E}.$




\textbf{Solution 4}

It's easier to calculate the number of lattice points inside a rectangle with vertices $(0,0)$, $(p,0)$, $(p,q)$, $(0,q)$. Those lattice points are divided by the diagonal $y = \frac{p}{q} \cdot x$ into $2$ halves. In this problem the number of lattice points on or below the diagonal and $x \ge 1$ is

 \begin{verbatim}
$\frac{1}{2} [(p+1)(q+1) - d] + d - (p+1)$, $d$ is the number of lattice points on the diagonal, $d \ge 2$
\end{verbatim}

$(p+1)(q+1)$ is the total number of lattice points inside the rectangle. Subtract the number of lattice points on the diagonal, divided by 2 is the number of lattice points below the diagonal, add the number of lattice points on the diagonal, and subtract the lattice points on the $x$ axis, then we get the total number of lattice points on or below the diagonal and $x \ge 1$.

 There are $900$ lattice points in total. $300$ is $\frac{1}{3}$ of $900$. The $x$ coordinate of the top-right vertex of the rectangle is $30$, $\frac{1}{2} \cdot 30 \cdot 20 = 300$. I guess the $y$ coordinate of the top-right vertex of the rectangle is $20$. Now I am going to verify that.  The slope of the diagonal is $\frac{20}{30} = \frac{2}{3}$, there are $11$ lattice points on the diagonal. Substitute $(p,q)=(30, 20)$, $d=11$ to the above formula:

 \begin{verbatim}
$\frac{1}{2} [(30+1)(10+1) - 11] + 11 - (30+1) = 300$
\end{verbatim}

Because there are $11$ lattice points on line $y = \frac{2}{3}x$, if $m < \frac{2}{3}$, then the number of lattice points on or below the line is less than $300$. So $m = \frac{2}{3}$ is the lower bound.

 Now I am going to calculate the upper bound. From $\frac{b}{a} < \frac{b+1}{a+1}$,

 $\frac{2}{3} = \frac{18}{27} < \frac{19}{28}$
$\frac{2}{3} = \frac{20}{30} < \frac{21}{31} < \frac{19}{28}$

                    \begin{figure}[H]
                        \centering
                        \includegraphics[width=0.5\textwidth]{Solution/AMC_12_12B_2021_25_Solution_Figure_2.png}
                    \end{figure}
                    

 If $m = \frac{21}{31}$, I will calculate by using the rectangle with blue vertex $(p,q) = (31, 21)$, then subtract lattice points on line $x = 31$, which is $21$. There are 2 lattice points on the diagonal, $d=2$.

 \begin{verbatim}
$\frac{1}{2} [(31+1)(21+1) - 2] + 2 - (31+1) - 21 = 300$, same as that of $m = \frac{2}{3}$
\end{verbatim}

If $m = \frac{19}{28}$, I will calculate by using the rectangle with red vertex $(p,q) = (28, 19)$, then add lattice points on line $x = 29$ and $x = 30$, which is $19 + 20 = 39$. There are 2 lattice points on the diagonal, $d=2$.

 \begin{verbatim}
$\frac{1}{2} [(28+1)(19+1) - 2] + 2 - (28+1) + 39 = 301$, $1$ more than that of $m = \frac{2}{3}$
\end{verbatim}

When $m$ increases, more lattice points falls below the line $y = mx$. Any value larger than $\frac{19}{28}$ has more than $301$ lattice points on or below $y = \frac{19}{28} x$. So the upper bound is $\frac{19}{28}$.

 $\frac{19}{28}-\frac{2}{3}=\frac{1}{84}$, $\boxed{\textbf{(E)} ~85}$.

 



\textbf{Solution 5}

The lower bound of $m$ is $\frac23 = \frac{20}{30}$. Inside the rectangle with vertices $(0,0)$, $(30,0)$, $(30,20)$, $(0, 20)$ and diagonal $y = \frac23 x$, there are $(30-1)(20-1) = 551$ lattice points inside, not including the edges. There are $9$ lattice points on diagonal $y = \frac23 x$ inside the rectangle, $551 + 9 = 560$. Half of the $560$ lattice points are below diagonal $y = \frac23 x$, $560 \cdot \frac12 = 280$. There are $20$ lattice points on edge $x = 30$, $280 + 20 = 300$. Once $m < \frac23$, the $9$ lattice points on diagonal $y = \frac23 x$ will be above the new diagonal, making the number of lattice points on and below the diagonal less than $300$.

 Now we are going to calculate the upper bound by the following formula:

 \begin{verbatim}
The number of lattice points inside rectangle $(0,0)$, $(a,0)$, $(a, b)$, $(0, b)$ and below diagonal $y = \frac{b}{a}x$ is $\frac{(a-1)(b-1)}{2}$, where $a$ and $b$ are relatively prime.There are $(a-1)(b-1)$ lattice points inside the rectangle. Because $a$ and $b$ are relatively prime, the only lattice points on the diagonal are $(0,0)$ and $(a,b)$. By symmetry, half of the lattice points are below the diagonal.
\end{verbatim}


                    \begin{figure}[H]
                        \centering
                        \includegraphics[width=0.5\textwidth]{Solution/AMC_12_12B_2021_25_Solution_Figure_3.png}
                    \end{figure}
                    
$\frac{2}{3} = \frac{20}{30} < \frac{21}{31} < \frac{19}{28}$

 When $a = 31$, $b = 21$, $\frac{(31-1)(21-1)}{2} = 300$.

 When $a = 28$, $b = 19$, $\frac{(28-1)(19-1)}{2} = 243$. Below the line $y = \frac{19}{28} x$, there are $19$ lattice points on line $x = 28$, $19$ lattice points on line $x = 29$, $20$ lattice points on line $x = 30$, $243 + 19 + 19 + 20 = 301$. 

 More lattice points fall below the line $y = mx$ as $m$ increases. There are more than $301$ lattice points on and below the line for any $m$ greater than $\frac{19}{28}$. Therefore, the upper bound is $\frac{19}{28}$.

 $\frac{19}{28}-\frac{2}{3}=\frac{1}{84}$, so $1+84=\boxed{\textbf{(E)} ~85}$.

 \textbf{Remark} $\lfloor \frac{b}{a} k \rfloor$ is the number of lattice points on line $x = k$, below line $y = \frac{b}{a} x$ and above the $x$ axis, where $k$ is an integer and $0<k<a$. Therefore, $\sum_{k=1}^{a-1} \lfloor \frac{b}{a} k \rfloor$ is the number of lattice points inside the rectangle $(0,0)$, $(a,0)$, $(a, b)$, $(0, b)$, below diagonal $y = \frac{b}{a} x$. If $a$ and $b$ are relatively prime, $\sum_{k=1}^{a-1} \lfloor \frac{b}{a} k \rfloor = \frac{(a-1)(b-1)}{2}$, as explained in solution 5. This problem is about finding the upper and lower bound of $\frac{b}{a}$, given $\sum_{k=1}^{30} \lfloor \frac{b}{a} k \rfloor = 300$. The same problem can have geometric representation as stated in the original problem, or algebraic representation as stated here.

 



\textbf{Solution 6}

As the procedure shown in the Solution 1, the lower bound of $m$ is $\frac{2}{3}$, whose equation is $2x - 3y = 0$.

 Now, we are going to find the upper bound of $m$.

 The signed distance between a point $(x_0,\ y_0)$ and the line $2x - 3y = 0$ is 
 \[d = \frac{2x_0 - 3y_0}{\sqrt{2^2 + 3^3}}\]

 If $(x_0,\ y_0)$ is the lowest lattice point above the line, then it should have the largest values of coordinates among all the solutions of $2x_0 - 3y_0 = 1$, because among the points with the shortest distance from the line, the slope of the farthest point to the origin is the smallest. Solve the inequality:
 \[2x - 3y = 1,\ 1 \leq x \leq 30\]

 We obtain $(x_0,\ y_0) = (28,\ 19)$, so the upper bound of $m = \frac{19}{28}$.

 $\frac{19}{28}-\frac{2}{3}=\frac{1}{84}$, $\boxed{\textbf{(E)} ~85}$.

 


\end{enumerate}